\documentclass[12pt]{article}

% Language setting
\usepackage[english]{babel}
    
% Set page size and margins
% Replace `letterpaper' with`a4paper' for UK/EU standard size
\usepackage[letterpaper,top=2cm,bottom=2cm,left=2.5cm,right=2.5cm]{geometry}

% Bib
%\usepackage[authoryear,round]{natbib}


\usepackage[T1]{fontenc}
\usepackage[utf8]{inputenc}


% This is used to find bad characters causing compilation errors
%\DeclareUnicodeCharacter{0301}{*************************************}


%For diagrams
\usepackage{tikz}
\usetikzlibrary{positioning,arrows.meta,calc}



% bibliography
\PassOptionsToPackage{natbib=true}{biblatex}
\usepackage{csquotes}
\usepackage[style=authoryear,maxcitenames=2, maxbibnames=99, uniquelist=false,  giveninits=true, doi=false, isbn=false, url=false, eprint=false, bibencoding=utf8]{biblatex}
\AtEveryBibitem{\clearlist{language}}
\addbibresource{code/documents/populist.bib}

%for comments in paper
\usepackage[status=draft,inline,index]{fixme}
%\usepackage[status=final,inline,index]{fixme}
\fxsetup{theme=color,margin=false,mode=multiuser}
\FXRegisterAuthor{rl}{rle}{RL}
\FXRegisterAuthor{ng}{nge}{NG}
\FXRegisterAuthor{od}{ode}{OD}
\FXRegisterAuthor{sk}{ske}{SK}

% Useful packages
\usepackage{amsmath}
\usepackage{graphicx}
\usepackage{booktabs, multicol, multirow, tabularx}
\usepackage{rotating}

% for figures, subfigures
\usepackage{caption}
\usepackage{subcaption}

\usepackage{booktabs}
\usepackage{multirow}
\usepackage[capposition=top]{floatrow}
\usepackage{float}


\usepackage{graphicx}


% for table
\usepackage{longtable}
\usepackage{makecell}
\usepackage{xcolor}
\usepackage{colortbl}

\usepackage{pdflscape} % rotate wide table

% This reduces the large amount of white space we had before and after each equation 
%\usepackage[nodisplayskipstretch]{setspace}
\expandafter\def\expandafter\normalsize\expandafter{%
    \normalsize%
    \setlength\abovedisplayskip{8pt}%
    \setlength\belowdisplayskip{6pt}%
    \setlength\abovedisplayshortskip{-8pt}%
    \setlength\belowdisplayshortskip{6pt}%
}

\usepackage{verbatim}
\usepackage{stmaryrd}
\usepackage{setspace}
\usepackage{layout}

% For Grammarly
\usepackage{endnotes}
\usepackage[colorlinks=true, allcolors=blue]{hyperref}
\newcommand{\beginQuotable}[1][quote name]{\phantomsection\label{Q:#1}}
\def\endQuotable{}

% Spacing
\setstretch{1.5}

%for cross-referencing
%\usepackage{xr}

%\doublespacing

\title{Decomposing the Rise of the Populist Radical Right\thanks{Danieli: Tel Aviv University and CEPR,  orendanieli@tauex.tau.ac.il. Gidron: Hebrew University of Jerusalem, noam.gidron@mail.huji.ac.il. Kikuchi: UCSD, shkikuchi@ucsd.edu. Levy: Tel Aviv University and CEPR, roeelevy@tauex.tau.ac.il. We are grateful to Daron Acemoglu, David Autor, Amir Beck, Leonardo Bursztyn, Nathan J. Canen, Eve Colson-Sihra, Natalia Emmanuel, Benjamin Enke, Makoto Fukumoto, Jeffry Frieden, Clément Herman, Alexey Makarin, Yotam Margalit, Vincent Pons, Giacomo A. M. Ponzetto, Moses Shayo, Carolyn Tsao, Ebonya Washington, and seminar participants at the ASSA meeting, the Tel Aviv University BMI Conference, Bocconi University, Bologna University, Computational Political Psychology Network, CEPR Political Economy Programme and CEPR Research and Policy Network on Populism Joint Symposium, EIEF, Hebrew University, Heinz College CMU, IASEI Annual Conference, the Israel Economic Association, IZA, Japan Applied Econometrics Conference, Japanese Society for Quantitative Political Science Winter Meeting, the Populism Seminar, NBER Political Economy, the Social Economics Conference, PEPES Seminar, Princeton, Stockholm Institute of Transition Economics, Swedish Institute for Social Research, Tel Aviv University, UBC Behavioral/PE Seminar, UBC Labor Seminar, UPF, and UZH for helpful comments and suggestions. We thank Maor Ankri, Noga Boneh, Yochai Cohen Or, Elad Guttman, Gilad Hurvitz, Amit Jacobson, Edan Kasan, Elior Mihelashvili, Liam Miller, Leanne Morical, Maya Orenstein, Wei Pin, Ori Shoham, Shlomit Shtein, Eran Yogev, and Ofir Zeira for excellent research assistance. Funding for this project was provided by ISF Grant 2319/21, the Foerder Institute for Economic Research, the Pinhas Sapir Center for Development, and the Archimedes Center at Tel Aviv University.}}

\author{Oren Danieli, Noam Gidron, Shinnosuke Kikuchi, Ro'ee Levy}


\date{\today} 
%\begin{normalsize} Click \href{https://www.dropbox.com/s/lpabpm9is2ei4s1/PopulistVoting\%20website.pdf?dl=0}{here} for latest version \end{normalsize} }


\begin{document}

\newcommand{\ivsWeights}{We use IVS sampling weights within countries and equally weight each country.}
\newcommand{\ivsWeightsCountryInverse}{We run this analysis in every country (using IVS sampling weights within countries), and aggregate over countries using the inverse of the share of PRRP support in the 2017–2020 wave (the same country weights as the main analysis).}
\newcommand{\PRRPscore}{The PRRP score measures potential support for PRRP based on voter opinions. It uses the fitted value from a LASSO regression predicting PRRP support using all voter opinions. See Section \ref{sec:data-attributes} for more details. }


% This fixed placement of floats, taken from: https://aty.sdsu.edu/bibliog/latex/floats.html
   \renewcommand{\topfraction}{0.9}	% max fraction of floats at top
    \renewcommand{\bottomfraction}{0.8}	% max fraction of floats at bottom
    %   Parameters for TEXT pages (not float pages):
    \setcounter{topnumber}{2}
    \setcounter{bottomnumber}{2}
    \setcounter{totalnumber}{4}     % 2 may work better
    \setcounter{dbltopnumber}{2}    % for 2-column pages
    \renewcommand{\dbltopfraction}{0.9}	% fit big float above 2-col. text
    \renewcommand{\textfraction}{0.07}	% allow minimal text w. figs
    %   Parameters for FLOAT pages (not text pages):
    \renewcommand{\floatpagefraction}{0.7}	% require fuller float pages
	% N.B.: floatpagefraction MUST be less than topfraction !!
    \renewcommand{\dblfloatpagefraction}{0.7}	% require fuller float page



\maketitle

% For Grammarly 
\def\grammarlyVersion{0}

\if\grammarlyVersion1
\let\footnote=\endnote 
\pagenumbering{gobble}
\righthyphenmin=62
\lefthyphenmin=62
\fi

% Remove abstract title
\renewcommand{\abstractname}{\vspace{-\baselineskip}}

 
% Abstract - closer to 200 words, we'll shorten before submissions, good enough for website and job market
\begin{abstract}
\setstretch{1.0}
\noindent 
Support for populist radical right parties in Europe has increased dramatically in recent years. We decompose the rise of these parties from 2005 to 2020 into four components: changes in party characteristics (positions and leadership demographics), changes in voter attributes (opinions and demographics), changes in voter priorities, and a residual. We merge wide datasets on party characteristics and voter attributes and estimate voter priorities using a probabilistic voting model. The primary driver behind the increased success of the populist radical right lies in voters' changing priorities. In particular, voters increasingly prioritize cultural issues over economic issues. This allows populist radical right parties to tap into a preexisting reservoir of culturally conservative voters, who tend to be older and without a college degree. By contrast, changes in voter attributes and party characteristics do not play a major role in the recent rise of populist radical right parties. We provide reduced-form evidence supporting these results. First, voters increasingly self-identify ideologically based on their cultural opinions rather than their economic opinions, demonstrating that their priorities have changed. Second, voters have not adopted more populist right-wing opinions, explaining why voter opinions do not drive the rise of the populist radical right.  Third, while parties' positions have changed, these changes are inconsistent with the main supply-side hypotheses for the increase in populist radical right support.



\end{abstract}

{\footnotesize \textbf{JEL codes:} D72; P00; H80.}
\thispagestyle{empty}


\clearpage{}

% View page layout to determine graph width
%\layout*

% Add tex values
\input{output/values/values}
\input{output/values/modelValues}
\input{output/values/manifestoValues}


\setcounter{page}{1}
\noindent
% Motivation I - RL: PERFECT! 
The rise of populist radical right parties (PRRP) is one of the most significant political developments in recent decades. Once on the fringes of European politics, PRRP now capture more than 20\% of the vote, as shown in Figure \ref{fig:populist_support}. PRRP are increasingly present not only in parliaments but also in governments, influencing policymaking on multiple issues, ranging from immigration to welfare \citep{akkerman2016radical, rathgeb2021study}. In certain countries, populist governments have resulted in lower GDP and significant erosion of democratic norms and institutions \citep{mccoy2019toward, funkePopulistLeadersEconomy2020}.


While a growing literature has identified various factors affecting PRRP support \citep{rodrik2018populism,fetzerDidAusterityCause2019,mudde2019far,tesei2022}, there is still no consensus on the primary factor behind the extensive rise of these parties \citep{guriev2022political}. Moreover, there is no agreement on the broad mechanisms behind the rise of PRRP, including whether this trend is shaped by supply or demand forces \citep{golder2016far, guisoDemandSupplyPopulism2017, dalboEconomicSocialOutsiders2020, hall2023political}. Scholars focusing on the supply side consider how shifts in party positions, such as the moderation of PRRP on certain issues, can explain the growing support for these parties \citep{akkerman2015gender}. On the demand side, an ongoing debate centers on two potential explanations. A common view argues that support for PRRP may have increased because public opinion has become more nativist or authoritarian \citep{hangartner2019does, ballard2021economic}. The alternative view argues that voters have changed their priorities rather than their opinions: nativist or authoritarian opinions have always existed, but they now shape vote choice more strongly \citep{bartels2017wave,bonikowski2017ethno, enke2022morals}.

% Summary Paragraph 
In this paper, we compare the explanatory power of these three classes of arguments: shifts in party characteristics (positions and leadership demographics), changes in voter attributes (opinions and demographics), and changes in voter priorities. We create a new dataset that links data on party characteristics with voter attributes, and estimate voter priorities using a probabilistic voting model. To directly compare the arguments, we develop a novel decomposition method for voting behavior, building on common methods for wage decomposition \citep{fortin2011decomposition}. This descriptive method shows that, of the three channels, rising support for PRRP is mostly driven by changing priorities. Specifically, we find that over the past two decades, voters have increasingly prioritized cultural issues over economic issues. 

% Model 
We start by outlining a multidimensional probabilistic voting model of how party characteristics, voter attributes, and voter priorities jointly determine which party each voter prefers. In the model, a voter's utility from supporting a specific party is a function of that party's characteristics (supply) weighted by the voter's individual voting weights (demand).\footnote{In Section \ref{sec:results}, we show that our results are robust to an alternative bliss-point model.} Each weight corresponds to a specific party characteristic and can be positive or negative, depending on whether the voter favors or dislikes this characteristic. 

The two demand channels, voter attributes and voter priorities, affect voting decisions through the voting weights. We allow the weights to vary across voters by assuming that they are a function of the voter attributes. The mapping between voter attributes and the weights is determined by the voter priorities, which correspond to the parameters of the utility function. We refer to these parameters as priorities since they determine which party characteristics voters prioritize given their attributes. This model structure allows us to separately evaluate the two categories of changes in demand discussed in the literature. The first is a change in attributes, reflecting a shift in voters' opinions or demographic composition (e.g., a shift toward more nativist attitudes). The second is a change in priorities, reflecting a shift in the salience of specific topics or the importance voters attach to them (e.g., nativist topics become more consequential in shaping vote choice).


% Data  
We measure party characteristics and voter attributes by creating a novel dataset that links the Comparative Manifesto Project (CMP), parties' leadership demographics, and the Integrated Values Survey (IVS). The CMP provides data on 56 party positions based on the parties' platforms. We supplement the party positions with data we compiled on the demographics of party leaders and the share of female legislators. The IVS includes a rich and consistent set of voter attributes, including their opinions and demographics. We analyze \numXVariables{} attributes over three waves: 2005--2009, 2011--2013, and 2017--2020.\footnote{While the IVS data provides rich information on voter attributes, it does not allow us to determine whether respondents would actually vote. Hence, we do not directly study turnout, despite its importance \citep{guisoDemandSupplyPopulism2017}. Reassuringly, as we discuss in Section \ref{sec:data-attributes}, the large rise in PRRP vote share in our survey data is similar to the actual rise in PRRP support in elections held around the same time.} 
Our merged dataset includes over 50,000 respondents in \countriesNum{} European countries. 

We use our merged dataset to estimate voter priorities, separately for each survey wave. Since this dataset includes a rich set of voter attributes and party characteristics, our parameter space is high-dimensional. Hence, we leverage machine-learning techniques to develop a computationally feasible method to estimate the voter priorities.

% Decomposition 
Combining the data on parties and voters with the estimated priorities, we decompose the rise in support for PRRP between 2005--2009 and 2017--2020 into four components: party characteristics, voter attributes, voter priorities, and a residual. We quantify the contribution of each component to the rise in PRRP support across Europe.
Specifically, we calculate partial-equilibrium counterfactual changes in PRRP support when only one component shifts over time while the other three are held fixed. For example, to estimate the voter priorities component, we calculate the counterfactual rise in PRRP support if party characteristics, voter attributes, and the residuals are fixed at their 2017--2020 values, and only voter priorities shift from their 2005--2009 to their 2017--2020 values. 

%decomposition results!! 
Aggregating results from all countries, changes in voter priorities account for \perPrioritiesAllTime{}  of the overall rise of PRRP. By contrast, changes in voter attributes and in party characteristics account for only \perVoterAllTime{} and \perPartiesAllTime{}, respectively. The residual accounts for the remaining \perResidualAllTime{} and
is driven mainly by the entry of new PRRP into the political system. These results remain robust across various specifications, models, datasets, decomposition orders, and subregions. 

To unpack the mechanisms driving our results, we perform additional analyses focusing on each of the components. We use our model to examine how the changes in the estimated voter priorities affect voting decisions and find that voters today are more likely to prioritize parties' cultural positions rather than their economic positions. 
While the prioritization of cultural positions is widespread, it increases PRRP support, especially among older voters and voters without a college degree, who tend to hold more conservative cultural opinions. 


We complement these model-based results with independent direct survey evidence on changing priorities. We use voters' left-right self-identification as an alternative outcome that is not included in our main estimation. Consistent with our model results, we find that voters increasingly position themselves on the left-right scale based on their cultural opinions rather than their economic opinions. Additionally, we find an increase in the correlation between voters' cultural opinions and the cultural positions of the parties they support. 

%Voters
We then inquire into our surprising finding that changes in voter attributes barely contribute to the rise of PRRP. We analyze reduced-form trends in voter opinions and find almost no shift toward opinions associated with PRRP, when aggregating across countries. For example, in contrast to multiple media accounts, attitudes have not moved en masse toward greater opposition to immigration. This supports the hypothesis that a reservoir of potential PRRP voters existed before the dramatic rise in PRRP support. 

The size of this preexisting reservoir varies substantially across countries and can explain some of the cross-country variation in PRRP support. We use a spatial decomposition method to predict the counterfactual support that PRRP would have received in different countries if voters in all countries were facing the same choice set of parties. We find that the cross-country variation in voter attributes can partially explain the variation in PRRP support. Our exercise also allows us to estimate the PRRP reservoir in countries where populist parties struggle to compete due to first-past-the-post systems. For example, we extend our analysis to the U.S. and show that it has a relatively large PRRP reservoir.  

% Supply
On the supply side, we find that changes in party characteristics are not consistent with the main supply-side hypotheses for rising PRRP support. For example, we do not find that PRRP are moderating their cultural positions. 
\beginQuotable[GE]
While we rule out several supply-side hypotheses, our results cannot fully rule out general-equilibrium effects, in which supply shocks contributed to the rise of PRRP by affecting voter priorities. 
\endQuotable{}

Finally, we show that the increase in the residual is driven by the entry of new PRRP into the political system. Using voter attributes and voter priorities from the 2005--2009 wave, our model predicts that PRRP that entered the political system in the 2010s should have already received support in 2005--2009. Nevertheless, in multiple countries, this did not materialize, despite noticeable attempts. We discuss potential explanations for PRRP belated entry, such as coordination failures. 

% Literature - introduction
This paper contributes to the literature on populism by directly comparing supply and demand theories regarding the rise of PRRP. Previous papers estimated the roles of party positions and voter opinions in specific countries \citep[e.g.,][]{krasa2014policy, longuet2024party, umansky2025business}. This paper complements the literature on the causal effects of different shocks on PRRP support, including technological change \citep{anelliWeWereRobots2019}, globalization \citep{rodrik2020does}, financial crises and austerity policies \citep{fetzerDidAusterityCause2019, guiso2021financial}, immigration \citep{dustmann2019refugee}, trade \citep{colantone2018trade,autorImportingPoliticalPolarization2020,dippelEffectTradeWorkers2020}, and new media technology \citep{guriev3GInternetConfidence2020, tesei2022}. While our descriptive analysis is not intended to identify causal effects, it points to the channels through which these shocks operate. For example, if trade shocks increased PRRP support, our results suggest these shocks changed voters' priorities rather than their opinions. Hence, our descriptive analysis highlights the importance of focusing on priorities as the outcome variable when studying PRRP support.

A smaller literature estimates voters' priorities and their behavior in different contexts \citep{johns2010,de2013dynamics,kendall2015voters,sides2019identity,le2023campaigns}. We contribute to this body of research by estimating a comprehensive set of priorities across Europe. Our analysis provides empirical support for recent papers discussing the prioritization of cultural issues \citep{kriesi2008west, enkeMoralValuesVoting2020,Bonomi2021, gethin2021brahmin, enke2022morals, margalit2022}. 
%Of course, comprehensively estimating priorities across individuals and over time comes at the cost that our estimates are not causal. 

Related literature shows that PRRP can gain support by activating a reservoir of culturally conservative voters. \citet{cantoniPersistenceActivationRightWing2020} show that the AfD gained votes by activating right-wing cultural voters. \citet{voigtlander2012persecution} find that the Nazis gained support by activating latent historical reservoirs of anti-Semitic attitudes. The importance of salience, legitimacy, or priorities in activating conservative cultural attitudes has also been shown in the U.S. \citep{bursztynExtremeMainstreamErosion2020}, Italy \citep{magistro2021changing}, the United Kingdom \citep{sobolewska2020brexitland}, Western Europe \citep{dennison2019rising}, as well as in comparative research \citep{bartels2023democracy,bonikowski2017ethno}. 
%\beginQuotable[testImmigrationQuote]Moreover, across Western Europe, PRRP gain electorally when and where voters attach greater importance to the issue of immigration \citep{dennison2019rising,dennison2020issue}.\endQuotable{} 
We directly compare voter priorities and voter opinions and show that changes in priorities are the main driver behind the recent rise of PRRP support across Europe. 

\beginQuotable[decompositionLit]
Finally, we contribute to the literature by applying a decomposition method to study trends in voting behavior. This descriptive method made notable contributions in other fields, such as the inequality literature \citep{oaxaca1973male,juhn1993wage, dinardo1996labor}, as it uncovers essential stylized facts, compares competing potential drivers, and focuses causal research on the most important research questions. While we focus on the rise of PRRP, this method can be used to study other political developments as well. 
\endQuotable{}



\section{Voting Model}\label{sec:model}

In this section, we develop a voting model that provides a parsimonious unifying framework for how party characteristics, voter attributes, and voter priorities codetermine vote choice. 

We assume a standard utility maximization framework where voters support the party that maximizes their utility. Voter $i$'s utility from voting for party $j$ is a function of the party's characteristics weighted by her individual voting weights. Specifically, we assume %the following functional form for voter utility:  
\begin{equation} \label{eq:utility}
U_{ij} = w_{i}' z_{j} + \zeta_j+ \varepsilon_{ij}. 
\end{equation}
The $L$-dimensional vector $z_{j}$ represents party $j$'s characteristics, and $w_{i}$ is an $L$-dimensional vector of the corresponding voting weights. Each individual weight $w_{i}^{l}$ represents the impact of the corresponding party characteristic $z_j^l$ (positions and leadership demographics) on voter~$i$'s utility. The sign of the weight is positive when the voter favors a characteristic (i.e., the voter's utility increases when she supports a party with this characteristic) and negative when she opposes it. The weight's magnitude (in absolute terms) measures how much the voter cares about this characteristic compared to other characteristics. We use $\zeta_{j}$ to capture the common residual utility from voting for party $j$---an unobserved party property that increases the utility from supporting this party among all voters.  
This residual includes unobserved party characteristics that affect voters' utility from supporting the party, such as the party leader's charisma or the party's historical reputation. 
Finally, $\varepsilon_{ij}$ represents all unobserved idiosyncratic factors that affect voting decisions. 

We assume that the voting weights are a linear function of voter attributes:
\begin{equation}  \label{eq:weights}
w_{i} = x_{i}'\Phi+\beta.
\end{equation}
The $M$-dimensional vector $x_{i}$ represents the observed attributes (demographics and opinions) of voter $i$. The $M\times L$ matrix $\Phi$ determines how each voter attribute affects the weights voters place on every party characteristic. $\Phi_{ml} > 0$ means that a larger value of voter attribute~$x^{m}$ generates larger support for party characteristic $z^{l}$. For example, if $x^{m}$ measures voter support for redistribution and $z^{l}$ measures party support for larger welfare spending, we would expect  $\Phi_{ml}$ to have a positive value. We standardize the distribution of $x_i$, and therefore the $L$-dimensional intercept vector $\beta$ represents the average weight of the full population. Appendix Figure \ref{fig:model_flow} summarizes the model and how the different inputs map into voters' utility. 

We define a vector $\delta$ that captures the utility gain from each party $j$ that is common across voters and depends on the party characteristics and the common residual,
\begin{equation}
\label{eq:delta_def}
\delta_j := \beta'z_j+\zeta_j.
\end{equation}

\noindent Combining Equations \ref{eq:utility}, \ref{eq:weights} and \ref{eq:delta_def}, we can rewrite the utility as  a function of the interaction between voter attributes and party characteristics,

\begin{equation}
U\left(x_{i},z_{j},\varepsilon_{ij}\right)=x_{i}'\Phi z_{j}+\delta_{j}+\varepsilon_{ij}\label{eq:x*phi*z}.
\end{equation}

Equation \ref{eq:x*phi*z} can capture the second-order approximation of any functional form. It nests the standard bliss-point utility functions where voters vote for the party closest to them ideologically as discussed in Appendix \ref{sec:bliss}.
It also allows a more complex utility function, where voters vote based on multiple dimensions and where demographics, such as education, can also affect vote choice \citep{piketty2020capital, abou2021brahmin,gidron2022many}. 

\beginQuotable[IIAExplanation]Allowing the utility to vary by voter attributes ensures that the model does not rely on the independence of irrelevant alternatives (IIA) assumption. Simpler discrete choice models impose IIA, which restricts substitution patterns across parties. This yields implausible predictions, such as PRRP entry drawing support from center-right and green-party voters in equal proportions. We relax this assumption by letting the utility depend on the attributes of the decision-maker \citep[as in][]{berry1995automobile}. In our context, the interaction between voter attributes and party characteristics implies that different voters have different preferences over parties, resulting in more realistic counterfactuals. For example, our model predicts that the entry of a PRRP attracts voters disproportionately from parties that are closer substitutes for potential PRRP voters (e.g., center-right parties), rather than from parties with opposing positions (e.g., Green parties).\endQuotable{}

To take this model to the data, we assume that the unobserved idiosyncratic shock $\varepsilon_{ij}$ has a type-I extreme value distribution (Gumbel). Together with Equation \ref{eq:x*phi*z}, this assumption allows us to write the conditional probability of voting for a party as
\begin{equation} 
\label{eq:P(Z|X)}
P(z_j|x_i)= \frac{\exp(x_i' \Phi z_j+\delta_j)}{\sum_{k \in \mathcal{J}_{c,t}} \exp(x_i' \Phi z_{k}+\delta_{k})},
\vspace{6pt}
\end{equation}
where $\mathcal{J}_{c,t}$ is the set of all parties that are on the ballot in country $c$ in period $t$. 

\section{Data}\label{sec:data}
In this section, we discuss our data on party characteristics and voter attributes.\footnote{Appendix \ref{sec:data-linking} explains how we link these datasets.} \beginQuotable[Europe]Throughout, we focus on Europe for the following reasons. First, most European countries have proportional representation systems, allowing us to measure PRRP support directly. Second, PRRP have gained substantial strength in Europe in recent years \citep{rodrik2018populism}.  Third, these parties have long defined themselves in opposition to similar political and economic developments, such as European integration. Moreover, Appendix \ref{sec:classify_right} shows that European voters of PRRP share ideological worldviews. Fourth, these parties have formed networks of cooperation in supranational institutions such as the European Parliament, demonstrating that these parties perceive themselves as connected \citep{mcdonnell2020international}. 
\endQuotable{}

\subsection{Party Characteristics} \label{subsec:party_data}

\subsubsection{Party Classification} \label{subsec:party_class}
PRRP have at least three common characteristics \citep{europe2007populist}. First, these parties are nativist; i.e., they consider minorities as a threat to the purity of the ``real people.'' Second, these parties are authoritarian; i.e., they hold a ``belief in a strictly ordered society, in which infringements on authority are to be punished severely'' \citep[p. 29]{mudde2019far}. Third, these parties are populist; i.e., their politics is predicated on a moral opposition between corrupt elites and the unified people. 

\beginQuotable[populistClassify]We determine whether a party is a populist radical right party according to the PopuList \citep{rooduijn2024populist, rooduijn2023populistdata}, which classifies European parties based on experts' judgment. We define PRRP as parties that were classified as populist radical right parties in this data for any year since 2005.\footnote{The PopuList rarely classifies a party as a PRRP for only some years. To avoid interpreting reclassifications as changes in PRRP support, we define a party as PRRP if the PopuList ever defined it as such.} \endQuotable{} Appendix Table \ref{tab:prrp_parties} provides the full list of PRRP in our data. We use the Comparative Manifesto Project (CMP) data to classify parties into other party families.  See Appendix \ref{subsubsec:classification} for details. 
\endQuotable{}

\subsubsection{Party Positions: Comparative Manifesto Project} \label{subsec:CMP}
\beginQuotable[cmp] We measure party positions on various issues using the CMP \citep{volkens2020manifesto}. This dataset covers the manifestos (platforms) of parties contesting lower-house elections. The CMP codes what share of the manifesto is dedicated to each topic, and for many topics, it distinguishes between positive and negative mentions. Previous work shows that the CMP measures are strongly correlated with those from the Chapel Hill Experts Survey \citep{adams2019problem}. CMP positions also align closely with survey responses of members of parliament \citep{guenther2025political}. Furthermore, they have been shown to predict the parties' subsequent policy choices \citep{knill2010parties}. Our analysis includes all 56 CMP main (three-digit) variables, which are listed in Appendix Table \ref{tab:manifesto_vars}.
\endQuotable{}

We aggregate economic and cultural positions using two indices created by the CMP.\footnote{The distinction between economic and cultural issues is imperfect, as not all political issues neatly fall into one of the two categories \citep{cramer2016politics,gidron2025european}. Still, distinguishing between economic and cultural issues is a common feature in electoral politics research \citep{kitschelt1994transformation, Margalit_2019} and is analytically useful for our purposes.} The economic index aggregates 19 party positions on the government intervention--free market scale. The cultural index (originally called the society index) aggregates 11 cultural positions on a progressive-conservative scale, including traditional morality, national way of life, multiculturalism, and democracy. Both indices are constructed by adding conservative positions and subtracting liberal positions such that a higher value reflects more support for free markets or conservative cultural values. Appendix Table \ref{tab:manifesto_vars} details the positions included in each index. We verify the robustness of our results using multiple alternative indices in Section \ref{subsec:priorities-model}. 


\subsubsection{Leadership Demographics} \label{subsec:leadership}

\beginQuotable[demographicsMain]We augment our data on party positions with data on party leaders' demographics. Leadership demographics may drive PRRP support, especially in cases where mainstream parties are not demographically representative of all voters \citep{dalboEconomicSocialOutsiders2020}. We collect and code demographic information for the leaders of \finalparty{} party-election observations in our data, including gender, years of schooling, age, marital status, number of divorces, minority status (a single indicator for immigrant, non-Christian, non-white, or non-heterosexual), and the average log wage in the leader's previous occupation. In addition, we compile data from three sources on the share of female legislators for each party \citep{weeks2023mannerparteien, weeks2022data, AdamsBrackenGidronHorneOBrienSenk2023, adams2022data, GlobalLegislatorDatabase2024}. Appendix \ref{demogrphics} provides additional details, and Appendix Table \ref{tab:leader_code_book} describes the definitions of all demographic variables.  %\endQuotable
\endQuotable{}

\subsection{Voter Attributes: Integrated Values Survey}\label{sec:data-attributes}

We measure voter attributes using the Integrated Values Survey (IVS). The IVS is composed of two large-scale cross-national repeated surveys: the World Values Survey (WVS) \citep{haerpfer2021wvs} and the European Values Survey (EVS) \citep{evs2021trend}. This dataset provides several advantages for our analysis. First, it includes information on a variety of voter attributes, including demographics, social values, opinions, and party support. Second, many of the questions in the IVS are consistently asked over time. Third, the data covers a broad range of countries.

We study PRRP support in the 2005--2009, 2011--2013, and 2017--2020 survey waves. We focus mainly on the first and last waves, analyzing data from \countriesNum{} European countries, as only a subset of seven countries participated in the 2011--2013 wave. We include all European countries that were surveyed in both the 2005--2009 and the 2017--2020 waves, and in which at least 70\% of the voters supported a party that can be matched to the CMP, and at least 50\% of voters supported parties that existed in the first and last waves: Austria, Denmark, Estonia, Finland, France, Germany, Greece, Hungary, Iceland, Italy, Lithuania, the Netherlands, Norway, Poland, Portugal, Slovenia, Spain, Sweden, Switzerland, and the UK.\footnote{Support for parties that appear only in later waves is mechanically assigned to the residual component (see Section \ref{sec:residual}). When such parties account for a majority of the vote share, we cannot meaningfully distinguish the roles of party characteristics, voter priorities, and the residual. Hence, we exclude these countries from the main analysis. Section \ref{sec:decomp_robustness} shows similar results under alternative exclusion thresholds.}
Overall, we analyze \IVSObvsNum{} observations across the three survey waves (see Appendix Table \ref{tab:ivs_data} for a summary by wave). We use sampling weights in all of our analyses and impute missing values using random forests. 

Figure \ref{fig:populist_support_map} presents the support that PRRP received in the 2017--2020 survey wave. The correlation between these values and the support PRRP received in the election closest to the survey is \prrpSupportCor{}. Furthermore, we find a correlation of \prrpSupChangeCor{} between the change in PRRP support in the surveys and the actual change in their support in the closest election.\footnote{Discrepancies between survey results and election outcomes may arise for various reasons, including changes that occurred between when the survey was conducted and when elections took place, and trends in voter turnout not directly captured by the survey. Most WVS and EVS surveys do not include a question directly measuring turnout. We may be partially capturing turnout indirectly since respondents who abstain from voting can report not supporting any party.} 
 
To estimate changes over time while providing our model with as much flexibility as possible, we include in our analysis nearly all variables that appear in all three survey waves. Appendix Table \ref{tab:ivs_vars} describes the \demOpNum{} opinion and demographic variables included in our data. Our outcome variable is respondents' preferred party, defined as the party they said they would vote for or the party that appealed to them most. Appendix \ref{subsec:data_processing} provides more details on the data processing and variables analyzed.



To characterize potential PRRP supporters based on their attributes, we construct a ``PRRP score'' that aggregates opinions associated with PRRP support. We first run a LASSO regression and predict support for PRRP in the 2017--2020 survey wave based on voter opinions. We then predict for each voter in each wave whether they would vote for a PRRP based on their opinions and define the standardized fitted value as their PRRP score.\footnote{We ensure that all the country indicators are taken into account in the LASSO regression by not allowing the model to penalize these variables. We do not use the country indicators when calculating the PRRP score in order to allow for cross-country comparisons.} Appendix Figure \ref{fig:populist_score_cor} shows that the voter attributes most strongly correlated with the PRRP score are \textit{jobs should prioritize natives}, \textit{homosexuality (not) justifiable}, and \textit{don't want immigrant neighbors}. 

\subsection{Summary Statistics} \label{subsec:summ_stat}

Panel~A of Table \ref{tab:sum_stat} presents summary statistics of party characteristics for PRRP compared to other parties. 
As expected, PRRP are more conservative, particularly on cultural issues. The table also presents the average manifesto shares devoted to the five positions that most strongly distinguish PRRP from non-PRRP. It demonstrates clear differences between them: PRRP dedicate a substantial and increasing share of their platforms to nationalistic topics (e.g., positive mentions of ``national way of life''), which other parties rarely mention. \beginQuotable[demographicsMain2]There are also differences in demographics: PRRP leaders tend to have less education and more divorces per leader, and there are fewer women among PRRP legislators. \endQuotable{}

Panel~B of Table \ref{tab:sum_stat} presents summary statistics for voters.
As expected, PRRP supporters are more likely to be male and less likely to live in urban areas \citep{gidron2020populism}. Over time, PRRP supporters hold more culturally conservative opinions. 

\section{Estimation}\label{sec:estimation}

We estimate the model parameters (voter priorities) using a two-step procedure. 
%To save space, we can spare the summary
%We first estimate the matrix $\Phi$ mapping voter attributes to voting weights and the vector $\delta$ of the common utility from each party. These parameters fully determine the likelihood of voting for each party (Equation \ref{eq:P(Z|X)}). We estimate them using a penalized MLE, separately for each IVS wave. We then estimate $\beta$, the average voting weights placed on each characteristic, and $\zeta$, the residuals, based on Equation \ref{eq:delta_def}, using the estimands for $\hat{\delta}$ from all three waves.

\subsection{First Step: Estimation of $\Phi$ and $\delta$} \label{subsec:first-stage}

\beginQuotable[FixParameters]
We assume that the model parameters change over time but are equal across countries. This assumption allows us to estimate multiple dimensions of changes in voter priorities over time by aggregating data across all European countries.\footnote{In Section \ref{sec:results}, we show that the main decomposition results also hold under weaker assumptions. First, we allow the model parameters to differ across lower- and higher-income countries. Second, we estimate the model parameters for each country using all other countries. The results remain similar in both cases.} 

Appendix Figure \ref{fig:party_correlation} assesses this assumption using a simple approximation of our model that predicts support for each party separately (see Appendix \ref{sec:classify_right} for details). We find that the parameters predicting support for a specific party in one country also predict support for a similar party in another country. For example, a model trained to predict support for UKIP in the UK also performs well when predicting support for AfD in Germany.
This result is consistent with voter priorities being similar across the continent. The figure also shows that PRRP supporters are similar across countries and are different from supporters of other populist parties. 

\endQuotable{}

Our parameter space is high-dimensional and, therefore, we rely on machine-learning techniques to avoid overfitting. Our dataset is large, as we prefer not to make any prior assumptions regarding which combination of variables is important for explaining party support. Instead, we use all available variables and let the data determine which variables are relevant.
As a result, the dimension of the matrix $\Phi$ is approximately 6,500 (the number of voter attributes multiplied by the number of party characteristics). To avoid overfitting, we constrain $\Phi$ by imposing $\left\Vert \Phi\right\Vert <c$ for some constant $c$. We follow recent econometric research \citep{athey2021matrix}, and use the nuclear norm as our matrix norm for two reasons. First, the nuclear norm is known to generate low-rank matrices, which are easier to interpret and imply that the voters decide which party to support based on relatively few dimensions \citep{kitschelt1994transformation,poole2001d, kriesi2008west}. Second, the nuclear norm generates a convex optimization problem that is computationally easier to solve. 

Specifically, we estimate $\Phi$ and $\delta$ using a penalized maximum likelihood estimator (MLE). We obtain the following maximization problem using the likelihood derived in Equation \ref{eq:P(Z|X)}:
\begin{equation} \label{eq:MLE}
\max_{\Phi,\delta}\mathcal{L}\left(\Phi,\delta\right)-\lambda\left\Vert \Phi\right\Vert =\max_{\Phi,\delta}\sum_{i}\log\frac{\exp\left[x_{i}\Phi z_{j\left(i\right)}+\delta_{j(i)}\right]}{\sum_{k\in \mathcal{J}_{c,t}}\exp\left[x_{i}\Phi z_{k}+\delta_{k}\right]}-\lambda\left\Vert \Phi\right\Vert. 
\vspace{4pt}
\end{equation}
We estimate the parameters separately for each IVS wave. We solve this maximization problem using proximal gradient descent \citep{hastie2019statistical} and choose the value of the penalization parameter $\lambda$ using cross-validation.

\subsection{Second Step: Estimation of $\beta$ and $\zeta$} \label{subsec:second-stage}

In the second step, we use the estimates of $\hat{\delta}$ from the first step to estimate $\beta$, the mean value for the weight placed on each characteristic, and $\zeta$, each party's residual. In this step, we combine information from all three waves. 
To estimate $\beta$ and $\zeta$ based on Equation \ref{eq:delta_def}, we take the difference between two consecutive survey waves and get the following equation:
\begin{equation} \label{eq:oaxaca}
    \Delta^{t+1}_t {\delta}_{j} =  
    \underset{\text{Voter Priorities}}{\underbrace{\Delta^{t+1}_t  \beta \bar{z_{j}}}} +  
    \underset{\text{Party Characteristics}}{\underbrace{\bar{\beta} \Delta^{t+1}_t z_j }}  + 
    \underset{\text{Residual}}{\underbrace{\Delta^{t+1}_t  \zeta_{j}}},
\vspace{4pt}
\end{equation}
where $\bar{\beta}$ and $\bar{z_j}$ represent the average values over the two periods. This equation decomposes the overall changes in $\delta$ ($\Delta \delta$) into the contribution of changes in voter priorities ($\Delta \beta$), party characteristics ($\Delta z$), and the residual ($\Delta \zeta$).\footnote{Voter attributes do not affect $\delta$ since $\delta$ is defined as the common utility all voters receive from a party.} The coefficients $\Delta^{t+1}_t  \beta $ and $\bar{\beta}$ are estimated by regressing $\Delta_t^{t+1} \hat{\delta}$ on $\bar{z_{j}}$ and $\Delta^{t+1}_t z_j$. Once $\Delta^{t+1}_t  \beta $ and $\bar{\beta}$ are estimated, we calculate the estimates for $\hat{\beta}_t, \hat{\beta}_{t+1}$ and substitute them into Equation \ref{eq:delta_def} to calculate $\hat{\zeta}$.  

Since the number of party characteristics is relatively large compared to the overall number of observations (the total number of parties in each wave), we make two additional assumptions to avoid overfitting. First, we assume a linear trend in $\beta$ for countries that appear in the 2011--2013 intermediate wave. This simplifying assumption allows us to use the information from the intermediate wave to estimate the main trend between our two major waves in 2005--2009 and 2017--2020 more precisely. 
This assumption implies that $\beta_{t+2}-\beta_{t+1} = \beta_{t+1}- \beta_t$, and therefore, the sum of Equation \ref{eq:oaxaca} for $\Delta^{t+1}_t$ and $\Delta^{t+2}_{t+1}$ is 
%the trick here is that you can add the middle z to the sum and it cancels out 
\begin{equation}
\label{eq:estimating_beta}
    \Delta^{t+2}_t {\delta}_{j} =  
    \Delta^{t+2}_t  \beta \bar{z_{j}} +  
    \bar{\beta} \Delta^{t+2}_t z_j   + 
    \Delta^{t+2}_t  \zeta_{j},
\end{equation}
where the averages ($\bar{\beta},\bar{z_j}$) are taken over all three periods. 

Second, we assume that the combinations of party characteristics that generate the largest differences in utility among voters are the same factors that determine the average utility across all voters.\footnote{This assumption provides a useful approximation as voters typically support parties based on contested issues. However, it would be violated if all voters had a homogeneous taste for certain party characteristics.} 
We use the singular value decomposition of the matrix $\Phi$ estimated in the first stage, $\Phi=U \Sigma V^T$, and calculate the first $k$ columns of matrix $V$, $[V]_k$. Intuitively, these are the $k$ linear combinations of party characteristics that generate the largest utility differences across voters. We then assume that $\beta$ is a linear combination of the columns of $[V]_k$.  Appendix \ref{sec:beta_restrcit} provides more details.  The value of $k$ represents the dimension of the average voting weights, i.e., the number of dimensions of party characteristics that explain why certain parties are more popular on average. We use $k=5$ for our main analysis to allow for richer explanations for average party differences while still avoiding overfitting. Section~\ref{sec:results} shows that other $k$ values yield similar results.

\subsection{Estimation Results: 2017--2020 Voting Weights}\label{subsec:voting_weights}

In this subsection, we explore the variations in voting weights in the 2017--2020 wave.
We calculate the weights each voter places on each party characteristic by combining the observed voter attributes with the estimated voter priorities using Equation \ref{eq:weights}. The weights are measured in units of standard deviation to utility units, defined as how a one standard deviation increase in this party characteristic would affect voter utility.\footnote{To provide some intuition for these units, in a two-party system, an increase of one utility unit is equivalent to a change from a 50/50 vote share to a 73/27 share.}
We also compute aggregated weights for the economic and cultural CMP indices. To do so, we take the average of the weights of all variables included in each index and flip the sign of the weights for variables that enter the index negatively.

Appendix Figure \ref{fig:phi_all_a} presents the voter attributes with the largest coefficients (in absolute value) in the cultural and economic indices.\footnote{The weight indices are linear functions of the voter attributes, where the coefficients are linear combinations of the columns of matrix $\Phi$ that correspond to the positions comprising the CMP indices.}  For each index, we show the values of the ten largest coefficients in the 2017--2020 wave. The patterns we uncover provide face validity to our analyses. For example, we find that, all else equal, individuals who express confidence in unions tend to reward parties with left-wing economic positions (i.e., they put a large negative weight on an index of right-wing economic positions). Moving to the cultural index, individuals who believe jobs should prioritize natives reward parties with right-wing cultural positions. Appendix Figure \ref{fig:phi_demo} analyzes the weights on two indices of leadership demographics. We find that older voters and men place higher weights on parties led by older white men. This demonstrates that our model can capture complex voter preferences. Even though the model has over 100 variables to choose from, it detects that these variables are important specifically for this index but not as important for the other indices.


Next, in Figure \ref{fig:weights_static} we examine how the distribution of weights varies by party family and find that PRRP voters place relatively large weights on conservative cultural issues. We plot the 25th, 50th, and 75th percentiles of the weight distribution separately for supporters of each party family in the 2017--2020 survey wave. The top two panels present results for the two CMP indices, and the subsequent panels present the weights on the positions with the largest variation in weights across party families. The weights PRRP voters place on the economic index are similar to those of mainstream right-wing voters. By contrast, PRRP voters care more about conservative cultural positions (e.g., national way of life) than mainstream voters do. Reassuringly, Figure~\ref{fig:weights_static} also shows that voters of other party families put more weight on positions at the core of their party's agenda. For example, socialist voters place large positive weights on support for labor groups.
%For example, green voters place large positive weights on environmental protection, while socialist voters place large positive weights on support for labor groups.

\section{Decomposition \label{subsec:Dynamic-Analysis-Decomposing}}
We now turn to the decomposition of the rise of PRRP. 
The statistic that we decompose is $S_{P}^{c,t}$, the share of PRRP supporters in country $c$ in period $t$. This share is defined as the sum of the probabilities that each voter supports a PRRP (marked with the event $\Pi$),%\footnote{Formally, $P(\Pi|x_i;(\Phi_{t},\beta_{t}),Z_t^c, \zeta_t^c)=P(\arg\max_{j} U_{ij} \in \mathcal{P}_{c,t})$, where $\mathcal{P}_{c,t}$ denotes the relevant set of PRRP.}

\begin{equation}
\label{eq:prrp_share}
    S^{c,t}_{P}:=\int P(\Pi|x_i;(\Phi_{t},\beta_{t}),Z_t^c, \zeta_t^c)f^{c}_{t}(x_i)\text{d} x_i. 
\end{equation}

The parameters $(\Phi_t,\beta_t)$ denote voter priorities, $Z_t^c=\{z_{j,t}\}_{j \in \mathcal{J}_{c,t}}$ is the matrix of observed party characteristics, $\zeta_t^c = \{ \zeta_{j,t}\}_{j \in \mathcal{J}_{c,t}}$ is the vector of residuals for all parties in the option set ($\mathcal{J}_{c,t}$), and $f_t^c$ is the density of voter attributes in country $c$ in period $t$. Using this notation, the change in PRRP support between periods $t$ and $t+1$ is
\[
\Delta_{t}^{t+1}S_{P}^c=\int P\left(\Pi|x_{i};(\Phi_{t+1},\beta_{t+1}),Z_{t+1}^c,\zeta_{t+1}^c\right)f_{t+1}^c\left(x_{i}\right)\text{d}x_{i}-\int P\left(\Pi|x_{i};(\Phi_{t},\beta_{t}),Z^c_{t},\zeta_t^c\right)f^c_{t}\left(x_{i}\right)\text{d}x_{i}.
\]

This parameterization allows us to decompose $\Delta_{t}^{t+1}S^c_{P}$ into the sum of four components: voter priorities, voter attributes, party characteristics, and a residual. Formally:
%\footnote{We estimate the model such that the share of support in each party is pinned to their value in the data. Therefore, $S{c,t}_P$ will also match its sample value exactly. } 
{\footnotesize
\begin{equation}
\begin{array}{ccc}
\Delta_{t}^{t+1}S_{P}^c & = & 

\underset{\text{Voter Priorities}}{\underbrace{\int P\left(\Pi|x_{i};(\Phi_{t+1},\beta_{t+1}),Z_{t+1}^c,\zeta_{t+1}^c\right)f^c_{t+1}\left(x_{i}\right)\text{d}x_{i}-\int P\left(\Pi|x_{i};(\Phi_{t},\beta_{t}),Z_{t+1}^c,\zeta_{t+1}^c\right)f^c_{t+1}\left(x_{i}\right)\text{d}x_{i}}}\\

 & + & 

\underset{\text{Voter Attributes}}{\underbrace{\int P\left(\Pi|x_{i};(\Phi_{t},\beta_{t}),Z_{t+1}^c,\zeta_{t+1}^c\right)f^c_{t+1}\left(x_{i}\right)\text{d}x_{i}-\int P\left(\Pi|x_{i};(\Phi_{t},\beta_{t}),Z_{t+1}^c,\zeta_{t+1}^c\right)f^c_{t}\left(x_{i}\right)\text{d}x_{i}}}\\

 & + & 

\underset{\text{Party Characteristics}}{\underbrace{\int P\left(\Pi|x_{i};(\Phi_{t},\beta_{t}),Z_{t+1}^c,\zeta_{t+1}^c\right)f^c_{t}\left(x_{i}\right)\text{d}x_{i}-\int P\left(\Pi|x_{i};(\Phi_{t},\beta_{t}),Z_{t}^c,\zeta_{t+1}^c\right)f^c_{t}\left(x_{i}\right)\text{d}x_{i}}}\\


 & + & 

\underset{\text{Residual}}{\underbrace{\int P\left(\Pi|x_{i};(\Phi_{t},\beta_{t}),Z_{t}^c,\zeta_{t+1}^c\right)f^c_{t}\left(x_{i}\right)\text{d}x_{i}-\int P\left(\Pi|x_{i};(\Phi_{t},\beta_{t}),Z_{t}^c,\zeta_t^c\right)f^c_{t}\left(x_{i}\right)\text{d}x_{i}}}
\end{array}\label{eq:decomposition}
\end{equation}
}

This mathematical identity states that the overall rise in PRRP support is the sum of the rise driven by each of the four factors. Each component represents a partial-equilibrium counterfactual increase in PRRP support when only one input changes, while the others are held fixed \citep{dinardo1996labor}. We simulate these counterfactuals by calculating the probability that each voter supports each party (Equation~\ref{eq:P(Z|X)}) and aggregating across voters.  

After decomposing the trends in PRRP support in each country separately, we aggregate the results over all countries. We scale the results in each country by the inverse of the share of PRRP support in the 2017--2020 wave. 
This scaling guarantees that the results are not driven predominantly by countries with very high levels of PRRP support (e.g., Hungary). 


Quantifying the contribution of party characteristics, voter attributes, and voter priorities relies on our ability to measure them properly in the data. In Section \ref{subsec:summ_stat}, we show that the CMP and IVS include variables on party positions and voter opinions that clearly distinguish both PRRP and their voters, respectively. In Section \ref{sec:reduced_form}, we show that we are also able to capture trends in these variables. Moreover, we show that the opinion variables explain a large share of the variation in PRRP support across countries. Still, additional party or voter characteristics may be absent from our data. The decomposition exercise may still capture some changes in these unobservables if they are correlated with observable variables. 

To accommodate parties' entry and exit, we include all parties that participated in either period $t$ or $t+1$. We then set the residual $\zeta_{j,t}$ to $- \infty$ if party $j$ does not participate or is too small to appear in our data in period $t$. This ensures that the predicted voting share for this party is set to zero. Therefore, changes in PRRP support that are related to entry and exit are attributed to the residual component.
In Appendix \ref{sec:residual_appendix}, we further decompose the residual component to calculate the share of the residual driven by entry. 

We estimate standard errors using 200 bootstrap replications, resampling voters with replacement within each survey wave and re-estimating our two-stage model in each replication. 

\subsection{Interpretation}
\beginQuotable[limitationDecomposition] The decomposition in Equation \ref{eq:decomposition} sheds light on the drivers behind the rise of the European populist radical right. Decomposition methods inherently follow a partial-equilibrium approach \citep{fortin2011decomposition}. 
%In our case, each component represents a partial-equilibrium counterfactual increase in PRRP support when only one input changes. 
As in other decomposition exercises, the partial-equilibrium counterfactuals allow us to account for the share of each component in driving the increase in PRRP support, but cannot isolate whether and how factors affect each other and do not account for strategic responses between the components \citep{canen2023decomposition}.\endQuotable{}
%For example, parties may adjust their positions following changes in mass public opinion \citep{adams2009moves}. Alternatively, changes in voter opinions may be triggered by parties' strategies \citep{barber2019does}. 
Still, learning which broad mechanisms drive the increase in PRRP support allows us to rule out a set of common explanations for this consequential political development. Moreover, it highlights which factors should be studied further to explain the rise of PRRP. The rest of this section elaborates on what mechanisms each component measures and how they relate to existing theories on the rise of PRRP.

\paragraph{Voter Priorities}
This demand component captures changes in the parameters $\Phi$ and $\beta$. These parameters determine whether voters support or oppose each party characteristic and how they prioritize the different party characteristics, given their opinions and demographics. 
%Changes in priorities may increase PRRP support even in the absence of shifts in voter attributes or party characteristics. Specifically, 
For example, \citet{bartels2023democracy} argues that European PRRP gained support by activating a preexisting reservoir of voters with culturally conservative attitudes and not because of a dramatic change in attitudes.%\footnote{Regarding support for right-wing populism in the United States, see \citet{bartels2020ethnic}.} 


\paragraph{Voter Attributes}
This demand component captures changes in $f$, the distribution of voter attributes $x_{i}$. This component is associated with the dominant image of the rise of the populist radical right as a political tsunami: a swift and powerful shift in public opinion toward the ethnonationalism and authoritarianism of these parties. For example,
\citet{hangartner2019does} show that greater exposure to refugees in Greece fueled opposition to immigration.%\footnote{Similar findings have been documented in Austria \citep{rudolph2021europe} and Norway \citep{nordo2021scope} although not in some other contexts \citep{cools2021local, schaub2021strangers}.} 

%\citet{ballard2021economic} argue that people who live in regions exposed to trade shocks adopted more authoritarian values---which, in turn, pushed voters toward PRRP and related causes (e.g., Brexit). In addition, changes in voter attributes include demographic changes, such as an increase in unemployment, that have been associated with anti-immigration attitudes and greater PRRP support \citep{alganEuropeanTrustCrisis2017,guisoDemandSupplyPopulism2017,dehdari2022economic}.\footnote{Other work questions this relationship between economic hardship and support for PRRP \citep{mutz2018status}.}

\paragraph{Party Characteristics}
This component captures changes in the supply of party characteristics $Z$. 
%We expect this component to be large if the rise of PRRP is driven mainly by supply changes. 
For example, PRRP may have moderated their cultural positions in order to appeal to mainstream voters \citep{akkerman2015gender,lancaster2020value}. %Alternatively, PRRP could have shifted their economic positions toward welfare chauvinism, understood as generous welfare benefits that exclude those who are deemed as unauthentic members of the nation \citep{schumacher2016mainstream}. 
%An alternative set of hypotheses focuses on position changes in other parties. \citet{berman2021causes} argues that the convergence of center-left and center-right parties on economic issues in the 1990s opened a space for the PRRP parties. 
Alternatively, PRRP could have gained support because the left turned toward progressive cultural positions and alienated some working-class voters \citep{goodhart2017road}.%\footnote{Former leader of the German Social Democrats, Sigmar Gabriel, argued that ``Winning over the hipsters in California cannot make up for losing the workers of the Rust Belt'' \cite[247]{abou2020electoral}.} More generally, the divergence of parties on cultural issues may have resulted in partisan realignment \citep{longuet2024party}.

\paragraph{Residual} 
The last component captures changes in the residual $\zeta$ of all parties. Changes in the residual could reflect several channels through which support for PRRP might increase. First, the residual could capture the entry and exit of parties from the political system, which we model as having $\zeta_{j,t} = - \infty$ for parties that do not appear in our data in period $t$. Second, the increase in the residual component could reflect a change in unobserved party characteristics of either the PRRP or their competitors. Third, the residual component can include an unexplained increase in PRRP support due to model misspecification.
%Changes in the residual could generate an increase in support for PRRP in several ways. First, through changes in unobserved party positions, either by the PRRP or their competitors. Second, through an increase in the valence of PRRP, for example, if these parties had more charismatic leaders in recent years. Third, through model misspecification. 

\subsection {Order of Components}\label{subsec:order}
The order of the four components in the decomposition can affect the results, as it determines whether to fix the other components at their levels at the start ($t$) or end ($t+1$) of the period. This order matters when the interactions among the components also contribute to the rise of PRRP. For example, a shift toward more nativist opinions would have a larger impact when voters with those opinions also started prioritizing cultural issues. 

We set the residual as the last component in all specifications. Hence, when calculating the counterfactuals for party characteristics, voter attributes, and voter priorities, the residual component is always fixed at its value at $t+1$ (as in Equation \ref{eq:decomposition}). This implies that we quantify the impact of changes in our three main components as if the new entrants had already participated in the elections in period $t$. Otherwise, the three main components would be mechanically zero for new entrants, as support cannot grow for a party that has not entered yet ($\zeta_{j,t}=-\infty$).

The order of the remaining components could vary, and therefore, we verify in Section~\ref{sec:decomp_robustness} that our results are robust to different orderings. In our main specification, we set party characteristics as the third component. We use manifestos from period $t+1$ when the party did not exist in period $t$.\footnote{There is no need to impute party characteristics for parties that did not exist in period $t+1$ (exits), as in all counterfactuals in which characteristics are set to their values at $t+1$, residuals are also set to their values at $t+1$, which is $-\infty$. Hence, the counterfactual support for exits is zero regardless of their characteristics.} Therefore, the party characteristics component only captures changes in the characteristics of parties that existed in both waves. Setting party characteristics as the third component guarantees that when calculating the counterfactuals for voter attributes and voter priorities, all party characteristics are taken from the same time period.%\footnote{Choosing a different decomposition order would fix party positions to their value at period $t$. Since we impute the positions for new entrants based on their value in period $t+1$, we prefer fixing other party positions to their value in the same time period and not mixing manifestos from period $t+1$ for new entrants and period $t$ for other parties.} 

We set voter priorities and voter attributes as the first and second components, respectively, in our main specification. As a result, the change driven by priorities can be interpreted as the counterfactual change in PRRP support when only priorities change, and all other components are fixed at their values in the same time period (2017--2020). 

\section{Decomposition Results}\label{sec:results}
Figure \ref{fig:decomposition} presents our key finding, the aggregated decomposition results. We focus on the change in PRRP support by fixing the initial value to 0\% and the final value to 100\%. For countries that are unobserved in the 2011--2013 wave, we impute their decomposition values using the average of their 2005--2009 and 2017--2020 values. 

Figure \ref{fig:decomposition} shows that voter priorities explain \perPrioritiesAllTime{} of the overall increase in PRRP support between 2005--2009 and 2017--2020. Party characteristics and voter attributes explain only  \perPartiesAllTime{} and \perVoterAllTime{}, respectively.\footnote{A negative value implies that PRRP support would have been expected to decrease during this period based only on the change in the component with the negative value.}  
The remainder of the increase is driven by the residual. 

Appendix Table \ref{tab:decomposition_table} shows the results by country. Each column presents the counterfactual increase driven by a specific component in raw units. While there is some variation across countries (e.g., voter attributes matter more in Hungary than in Sweden), in most countries, voter priorities are more important drivers of PRRP support than voter attributes or party characteristics.

%The main takeaways from this figure is that voters' opinions and demographics did not change dramatically in a way that is consistent with supporting PRRP. Furthermore, any changes in party position either did not lead to more PRRP support or were offset by other changes that decreased their support. Finally, we learn that priorities not only changed in the last 20 years but that changing priorities can lead to dramatic shifts in voting behavior. 

%In Section \ref{sec:reduced_form} we show that these results are driven by a preexisting reservoir of potential populist voters who began to support PRRP as cultural issues gained importance. 
%On aggregate, voters' cultural opinions have remained similar, suggesting that a pool of culturally conservative voters existed before the rise of PRRP. Party positions shifted, but not in ways that led to increased support for PRRP. What has changed is that conservative voters increasingly support parties with positions that are more associated with their cultural beliefs. 

\subsection{Robustness} \label{sec:decomp_robustness}
In Figure \ref{fig:decomposition_robustness} and Appendix Figure \ref{fig:decomposition_robustness_app}, we present a series of analyses to confirm the robustness of our decomposition results. 

\beginQuotable[decOrder]\paragraph{Component Order} Following the discussion in Section \ref{subsec:Dynamic-Analysis-Decomposing},  we examine the degree to which our results depend on the order of the decomposition by alternating the order of party characteristics, voter attributes, and priorities. This generates $3!=6$ possible decomposition orders.\footnote{The residual component, which includes party entry, is always placed last in the ordering; otherwise, any component following it would be mechanically zero for new entrants (see Section \ref{subsec:order}).} We follow \citet{guriev2023curtailing} and present in Panel~A of Figure \ref{fig:decomposition_robustness} the average value of the six different orders, which is the Shorrocks-Shapley value of the decomposition \citep{shorrocks2013decomposition}. The results remain very similar.  \endQuotable{}


\paragraph{Leave-country-out} One possible concern is that despite our use of penalized MLE, we overfit the model parameters and therefore wrongly conclude that changes in the overfitted voter priorities explain a large share of the rise of PRRP. An additional concern is that using the same parameters for all countries could yield misleading results.  To address both of these concerns, we perform a leave-country-out exercise. For each country, we estimate all model parameters without data on that country ($\hat{\Phi}_{-c},\hat{\beta}_{-c}$). Since the voter priorities for each country are estimated using data only from other countries, overfitting is not possible.\footnote{We calculate the residual vector ($\hat{\zeta}_{-c}$) for each party in that country so that the predicted voting shares under these parameters fit the data. We then decompose the rise of PRRP in that country under these estimated parameters using Equation  \ref{eq:decomposition}. We aggregate the results from all countries as in our main estimation.} Panel~B of Figure \ref{fig:decomposition_robustness} shows that the results of the leave-country-out decomposition resemble those reported in our main analysis. The priorities component is somewhat smaller, but remains dominant.
Hence, changes in priorities that were estimated in other countries could predict a large share of the increase in support for PRRP in the remaining country. %The residual component is larger in this exercise (\leaveCountryOutResidualNum ) compared to our baseline (\perResidualAllTime), as the model fits the data better within sample.  

\paragraph{Bliss-Point Model} Panel C of Figure \ref{fig:decomposition_robustness} shows that our results remain qualitatively similar when we replace our baseline model with a bliss-point model. Appendix \ref{sec:bliss} describes a model where voters have an optimal party characteristic (bliss point) and vote for the party with characteristics closest to their bliss point. As explained in the Appendix, this model is similar to the model we describe in Section \ref{sec:model}.  

\beginQuotable[Cohorts]\paragraph{Cohorts} Our results suggest that voters' attributes, including age, are not driving PRRP support. 
%This is not surprising, as the age distribution has not changed substantially during this period. 
An alternative possibility is that a change in cohort composition, rather than aging, drives some of the rise in support for PRRP (see \citet{schafer2022cultural} for a discussion). In Panel D, we replace the age variable with a cohort variable (we cannot include both variables due to collinearity). While priorities remain more important than party characteristics and voter attributes,  the attributes component increases from \perVoterAllTime{} to \cohortAttributeNum{} when using cohorts. One potential explanation is that cohorts born earlier are more familiar with the historical dangers of PRRP, and thus less likely to vote for these parties, conditional on their opinions. \endQuotable{}


\paragraph{Party Rhetoric} \beginQuotable[styleMain]
Voters may choose their preferred party not only based on its positions or leadership demographics, but also on other characteristics such as rhetoric  \citep{dai2022politicians}. Therefore, we complement our measure of party characteristics with two aspects of rhetoric---complexity and insult score---based on the manifesto texts. We follow frontier methods in measuring these characteristics \citep[e.g.,][]{ma14ller2022effect,di2023keep}. Specifically, we measure complexity using the following three metrics.  First, we use lexical diversity (the moving average type-token ratio) to measure conceptual complexity. \citet{di2023keep} show that simpler texts, such as France’s ``easy to read and understand'' manifestos, score lower on this measure than official manifestos. Second, we use the average number of words per clause to measure syntactic complexity. Third, we use total word count to measure document-level complexity, as parties may shorten their manifestos to market simple messages. For example, Trump dramatically shortened the Republican Party platform in the 2024 election \citep{trumpPlatform24}. We measure the insult score as the average predicted share of people who would consider a clause insulting for the top 5\% most insulting clauses in each manifesto, based on the Perspective API \citep{lees2022new}. Appendix \ref{subsec:style} provides more details, and Appendix Table \ref{tab:manifesto_rhetoric_sum_stats} shows that PRRP use more insulting rhetoric than non-PRRP.
%\footnote{In contrast to party positions and leadership demographics, these measures of party rhetoric are uncorrelated with voters' attributes or party families. Hence, we do not to include them in our main analysis.} 
\endQuotable
Panel~E of Figure~\ref{fig:decomposition_robustness} shows that running our model with the additional party characteristics does not substantially change our results. 


\paragraph{Party Pledges} \beginQuotable[pledgesMain] One potential concern with the manifesto data is that it includes both declarative statements and concrete pledges to promote specific policies. Some voters may care more about the pledges, especially since parties have been shown to at least partially fulfill them \citep{thomson2017fulfillment}. Therefore, we use a large language model to classify clauses in party manifestos as either pledges or declarative statements, and rerun our analysis when party positions are defined based only on pledges. See Appendix \ref{subsec:pledges} for details. Panel~F of Figure \ref{fig:decomposition_robustness} shows that some of our results change in this exercise (the residual component increases and the party characteristics component decreases), perhaps because of the noise introduced when analyzing only subsets of the manifestos.\footnote{Interestingly, we find a weaker correlation between voters' opinions and parties' positions when positions are based solely on policy pledges, compared to declarative statements. Hence, the large residual may also be due to voters caring less about concrete policy pledges and more about expressive declarative language.} However, the voter priorities component remains more important than voter attributes or party characteristics. 
\endQuotable

\paragraph{Higher- and Lower-Income Countries}


\beginQuotable[highLowIncome] Our decomposition results remain consistent when estimating voter priorities separately for higher- and lower-income countries. In the main analysis, we assume that the voter priorities are the same for all voters, and in Appendix~\ref{sec:classify_right}, we show that this assumption is a reasonable approximation of the data. In Panels G and H of Figure \ref{fig:decomposition_robustness}, we relax this assumption and allow distinct voter priorities for countries above and below a GDP per capita of \$50,000 in 2021.
Higher-income countries are from Western and Northern Europe, and lower-income countries are from Eastern and Southern Europe. In both regions, voter priorities are the main driver of the rise of PRRP. However, the voter attributes component slightly increases PRRP support in lower-income countries and has the opposite sign in higher-income countries. This could be explained by a shift toward more culturally conservative opinions in some lower-income countries, particularly Hungary, as we discuss in Section \ref{subsec:attributes}.
\endQuotable{}


\paragraph{Additional Robustness Tests} 
In Appendix Figure \ref{fig:decomposition_robustness_app}, we conduct a series of additional robustness tests. First, we run our model using only voter opinions and party positions, and without the demographics of voters and parties. This increases the residual component, potentially due to less predictive power, but priorities remain much more important than party characteristics and voter attributes. Second, we show that our results are robust to different choices of the parameter $k$ (which sets the number of dimensions we use in the second stage of the estimation). Third, we show that our results are robust to the choice of countries. In our main analysis, we only include countries where at least half of voters support parties that existed in the previous survey wave. We can omit this rule and add the Czech Republic and Slovakia, or alternatively, make it stricter and remove Italy and Greece. In both cases, priorities remain substantially more important than party characteristics and voter attributes. 


\section{Detailed Evidence by Component} \label{sec:reduced_form}

In this section, we analyze the change in each component. We show that voter attributes are stable on aggregate, and party characteristics are not shifting in ways that are expected to increase support for PRRP. These findings strongly support the hypothesis that the rise in PRRP support is primarily driven by the ``activation'' of a preexisting reservoir of potential PRRP voters rather than a dramatic change in attitudes \citep{bartels2023democracy}. We begin by analyzing the most dominant component, voter priorities, showing that voters prioritize cultural issues over economic issues.


\subsection{Changes in Voter Priorities} \label{subsec:priority_results}

\subsubsection{Which Priorities Changed? }  \label{subsec:which_priorities}

\paragraph{Model-Based Evidence}\label{subsec:priorities-model}

Figure \ref{fig:weights_dynamic} shows that since the early 2000s, voters have placed more weight on cultural issues relative to economic issues. We isolate changes in voter priorities by fixing voter attributes at their values in the 2017--2020 survey wave. We then analyze the counterfactual changes in weights that are driven only by changes in voter priorities. We aggregate weights into two indices based on the CMP party position indices as before. Figure~\ref{fig:weights_dynamic} shows that the distribution of the weights placed on the CMP economic index became substantially more concentrated around zero in 2017--2020, compared to the weights placed on the CMP cultural index. This result provides cross-country empirical evidence for the claim that cultural issues have become more important than economic issues.

Appendix Figure \ref{fig:alternative_cmp} shows that the prioritization of cultural issues, compared to economic issues, does not depend on the choice of particular party position indices. We repeat the exercise in Figure \ref{fig:weights_dynamic} using two alternative party position indices proposed by \citet{bakker2013measuring} and \citet{prosser2014building}. The growing importance of cultural issues is robust to using these alternative indices.

\beginQuotable[moreIndices]Appendix Figure \ref{fig:alternative_indices} repeats the analysis of Figure \ref{fig:weights_dynamic} for three additional indices. In the left panel, we plot a foreign policy Hawk-Dove index based on \citet{kiratli2021hawkdove}. In the middle panel, we plot the European Integration Index proposed by the CMP. In the right panel, we plot an index for government efficiency, which we constructed based on the Governmental and Administrative Efficiency and Political Corruption positions. The CMP and the alternative economic and cultural indices, together with these three indices, cover variables that, on average, comprise \percindex{} of the full manifestos. Appendix Table ~\ref{tab:manifesto_vars} lists the variables composing each index. 

We find that the weights placed on all three indices have become more concentrated around zero, in contrast to the cultural index. This suggests that cultural issues receive a higher priority not only relative to economic issues but also relative to other measured issues. 
\endQuotable{}

\paragraph{Reduced-Form Evidence} \label{subsec:priorities-reduced}

We next provide complementary reduced-form evidence showing that voters prioritize cultural issues over economic issues.

\beginQuotable[caseStudy]
As an illustrative case study, we first examine the changes in priorities in Germany. We analyze the issues voters care about the most, based on the Politbarometer \citep{fgw2025politbarometer}, a long-running representative German political survey. Appendix Figure \ref{fig:germany_case_study} shows that in 2005, only a small share of German survey respondents considered cultural issues the most important political topic. Economic issues were considered much more important. Indeed, early elections that year were called because of the country's economic challenges, with disagreements over economic reforms at the core of the disputes between the main rival parties.\footnote{The word ``immigration'' was not mentioned even once in an academic report describing the political dynamics of these elections \citep{hough2006german}.} A decade later, the German political landscape had been completely transformed, with cultural issues, and specifically immigration, becoming far more important, as shown in the figure. This change is often attributed to Germany opening its borders to almost one million refugees in 2015--2016 \citep{schmitt2022changing, faas2019more}. It is amid this wave of increased salience of immigration that the populist radical right Alternative for Germany party not only surfed into parliament for the first time but also emerged as the third largest party. %While their importance decreased somewhat after 2016, immigration and other cultural issues remained more important than economic issues at least until 2019. 
Interestingly, based on the IVS, German voters became more sympathetic toward immigrants. This case study demonstrates that large exogenous shocks can shift voting behavior by changing priorities even as opinions move in the opposite direction. 
\endQuotable{}

However, scholars also recognize the limitations of using direct survey questions to estimate voters' priorities.  As \citet{sides2022bitter} note, ``it can be challenging to measure political priorities [...] People who rate an issue as more important, for example, do not appear to rely more on that issue when they choose between candidates.'' Experimental evidence from the United States shows that issue importance has little, if any, effect on actual vote choice \citep{leeper2020more}. Finally, survey questions asking which issues respondents prioritized \emph{at the ballot} are not consistently included over time or across countries.   

For these reasons, we employ an alternative reduced-form method to measure changes in priorities using respondents' left-right self-identification. We analyze responses to an IVS question that asks participants where they place themselves on a left-right political scale ranging from 1 (left) to 10 (right). This question forces voters to prioritize across issues, as they reduce their opinions on a variety of topics to a single dimension of ideology. 

To link left-right self-identification to opinions, we construct indices of cultural and economic opinions based on IVS questions. We classify IVS variables as cultural or economic variables if they resemble one of the variables included in the cultural and economic CMP indices, respectively (the variables are specified in Appendix Table \ref{tab:ivs_vars}). We then define the loading of each variable using principal component analysis in order to capture the common variation among the variables included in the index. 

Using these opinion indices, we analyze the relationship between economic and cultural attitudes and left-right self-identification over time \citep{de2013dynamics}. We first nonparametrically regress left-right self-identification on the economic and cultural opinion indices, separately for the 2005--2009 and 2017--2020 survey waves. Formally, writing $y$ as the left-right self-identification and $x^C, x^E$ as the cultural and economic opinion indices, we estimate 
$$E[y|x^C, x^E] = \hat{f}_t(x^C,x^E).$$

The function $f_t$ is estimated separately for each survey wave with nonparametric regressions. Using the estimation results, we predict the reported ideology for each combination of quantiles of the economic and cultural opinion indices (the quantiles are calculated based on the pooled 2005--2009 and 2017--2020 surveys). Finally, for each combination of the $x^C, x^E$ quantiles, we subtract the 2005--2009 predictions from the 2017--2020 predictions to examine how left-right self-identification changed, conditional on opinions
$$\Delta\hat{f}(x^C,x^E)=\hat{f}_{t+1}(x^C,x^E)-\hat{f}_t(x^C,x^E).$$

Figure \ref{fig:heat_map} shows that between the 2005--2009 and the 2017--2020 waves, the cultural dimension became more dominant in shaping left-right self-identification. Darker colors mean that individuals in specific quantiles of the economic and cultural indices in 2017--2020 are expected to identify as more right-leaning compared to individuals in similar quantiles in 2005--2009. We find that changes in left-right self-identification are primarily associated with voters' cultural opinions. Voters holding more conservative cultural opinions (the top of the figure) now define themselves as more right-leaning, and voters holding more progressive cultural opinions (the bottom of the figure) now define themselves as more left-leaning. By contrast, the economic index is very weakly associated with changes in left-right ideology. 

The increased importance of cultural issues is most salient for voters who have conflicting views on culture and economics (the top-left and bottom-right corners of Figure \ref{fig:heat_map}). Over time, these individuals' left-right self-identification became more closely associated with their cultural views than with their economic views. Considering that the dependent variable, left-right self-identification, is not used in our model estimation, this figure provides independent evidence that voters are increasingly prioritizing cultural issues over economic issues. 

As an alternative reduced-form measure of the growing importance of cultural issues in voting decisions, we use vote choice directly. Panel~A of Appendix Figure~\ref{fig:opinion_position_cor} shows the correlation between voter opinions and the positions of their preferred parties, separately for economic and cultural issues. We use the same voter opinion indices for economic and cultural issues as in the previous figure, and the same CMP indices for party positions as before. All indices are standardized at the country-wave level. The correlation between the cultural opinions of voters and the cultural positions of parties has increased over time. By contrast, the correlation between the economic opinions of voters and the economic positions of parties remained stable. This figure provides intuition for the kind of variation our model uses when estimating $\Phi$.\footnote{The estimated parameters solve the first-order conditions of Equation \ref{eq:MLE}, which imply that the observed correlations between the attributes of voters and the positions of the parties they support (as reported in Appendix Figure \ref{fig:opinion_position_cor}) are close to the correlations predicted by the model.} 

\paragraph{Discussion}
\beginQuotable[whichPrioritiesDiscuss]
The evidence that cultural issues have become relatively more important is consistent with several theoretical perspectives. First, \citet{inglehartTrumpBrexitRise2016} argue that the long-run shift toward progressive cultural values since the 1970s generated a cultural backlash among older and more traditional voters, thereby increasing support for PRRP through cultural grievance (but see \citet{schafer2022cultural} for a critical perspective). Second, \citet{enke2022morals} argue that values behave like a “luxury good”---rising income leads voters to place greater weight on non-economic issues, including cultural issues. Third, \citet{Bonomi2021, gennaioli2025presidential}, building on \citet{shayo2009model}, argue that shocks, such as immigration, can raise the salience of cultural conflict and induce voters to shift from class-based to cultural identities. They argue that even economic shocks, such as the China shock \citep{autorImportingPoliticalPolarization2020}, may heighten cultural cleavages when they differentially affect voters with different cultural values. Fourth, political entrepreneurs may successfully politicize cultural issues and thus increase their relative priority \citep{bornschier2013new, de2021politicizing, kashner2024preempting}. Finally, voters may prioritize cultural issues not because these issues have become more important but because the stigma associated with supporting parties with extreme cultural positions has declined \citep{valentim2025normalization}.\footnote{Appendix~\ref{sec:legitimacy} further examines the last hypothesis of declining stigma. We examine the trend in political opinions and analyze non-respondents, and do not find support for this hypothesis.} By contrast, the evidence is less consistent with theories that focus on the de-prioritization of economic issues, as we find that cultural issues also became more important than non-economic issues.\footnote{For example, our results are not consistent with voters de-prioritizing economic issues because more economic decisions are made by non-legislative actors, such as the ECB \citep{ezrow2014responding}.} 

\endQuotable{}


\subsubsection{For which Voters did Priorities Change?}\label{subsec:how}
\beginQuotable[howIntro]


We next study heterogeneity in changes in voters' priorities. We first examine which groups exhibited the largest increase in support for PRRP due to changes in priorities. In our main decomposition, we calculate the counterfactual rise in support for PRRP when only priorities change, while all other components are held fixed at their values in 2017--2020 (the first line of Equation \ref{eq:decomposition}). We can calculate the same counterfactual rise in PRRP support for any subgroup of the population (e.g., older voters) by aggregating only over that subgroup.\footnote{We calculate the rise in support for every subgroup in every country, and aggregate over countries using the same weights as the main analysis.} We then compare the rise in support for PRRP in every subgroup, relative to the increased support among the full population. 

Figure \ref{fig:hetero_prior_prrp} shows that the change in priorities increased PRRP support, particularly for older voters, men, voters without a college degree, and voters with lower income. The differences across these demographics are substantial. For example, the increase in PRRP support driven by priorities for voters above fifty years old is almost 40\% larger than that of the general population. We also directly examine whether the change in priorities increases PRRP support primarily among a reservoir of potential PRRP voters. Figure~\ref{fig:hetero_prior_prrp} shows that the rise in the priority component is driven entirely by potential PRRP supporters, those with an above-median PRRP score. The priority-driven increase in PRRP support for this segment is more than twice as large as in the general population. 

Why did the change in priorities increase support for PRRP much more substantially among these particular subgroups? One possibility is that some voters experienced a larger shift in priorities towards cultural issues. An alternative explanation is that a similar shift in priorities could mainly increase support among the reservoir of potential PRRP voters. Prioritizing cultural issues would increase PRRP support for voters with conservative cultural opinions, but not for voters with progressive cultural views.

To distinguish between these two channels, we examine which voters increasingly prioritized cultural over economic issues. 
Appendix Figure~\ref{fig:hetero_prior_change} presents the average change in the cultural/economic weights ratio for different subgroups, where a higher value indicates greater prioritization of cultural issues.\footnote{Formally, for each voter $i$ in the 2017--2020 survey wave, we calculate the log ratio of the absolute values of the weights on the cultural and economic indices, when using either priorities from the 2005--2009 wave or the 2017--2020 wave, holding voter attributes fixed at their 2017--2020 levels: 
\begin{equation}\label{eq:chane_ratio}
\log\frac{|w_{culture}^{2017-2020}(x_i)|}{|w_{econ}^{2017-2020}(x_i)|}-\log\frac{|w_{culture}^{2005-2009}(x_i)|}{|w_{econ}^{2005-2009}(x_i)|}    
\end{equation}
We then calculate the average across all voters in each demographic group.
} 
 We find that across all subgroups, voters prioritize cultural issues at the expense of economic issues over time, with small differences across groups. The largest changes appear among men and union members. Older voters shift only slightly more than younger voters, and voters with a college degree or higher income shift about as much as others. These patterns differ from those documented in Figure \ref{fig:hetero_prior_prrp}, suggesting that the heterogeneity documented in Figure~\ref{fig:hetero_prior_prrp} is not driven by heterogeneity in the change in priorities.

We next examine differences in cultural opinions across demographic groups. Appendix Figure \ref{fig:demo_prrp_score_bar} reveals substantial variation in the average PRRP score, consistent with the heterogeneity in Figure~\ref{fig:hetero_prior_prrp}. Older voters, voters without a college degree, voters with lower incomes, and non-unionized voters are all substantially more culturally conservative than others. 
These results are consistent with the second explanation. The prioritization of cultural issues is broad and occurs across the entire population. However, this shift translated into greater support for PRRP, mainly among a preexisting reservoir of voters with conservative cultural views. This reservoir is disproportionately older, less educated, and has a lower income. 


\paragraph{Discussion}
The differences we find across demographics are consistent with some of the prominent theories discussed in the last subsection. Finding that the change in priorities drives PRRP support among older voters and among voters with a higher PRRP score (Figure~\ref{fig:hetero_prior_prrp}) is consistent with the cultural backlash theory \citep{inglehartTrumpBrexitRise2016}, while the increased support among rural and less-educated voters---groups that suffered more from globalization---is consistent with \citet{Bonomi2021}. Finally, higher-income voters prioritize cultural issues at least as much as lower-income voters (Appendix Figure~\ref{fig:hetero_prior_change}), which is consistent with cultural issues becoming luxury goods \citep{enke2022morals}. 
\endQuotable{}

\subsubsection{When did Priorities Change?} \label{subsec:when}
\beginQuotable[timing] 
We extend our reduced-form analysis in Section \ref{subsec:which_priorities} to earlier IVS waves and find that the shift in priorities we document appears to be a post-2009 development rather than a smooth continuation of earlier trends. First, Panel~A of Appendix Figure \ref{fig:heatmap_1_4_subset} replicates the analysis of Figure \ref{fig:heat_map}, examining how left-right self-identification changed for voters depending on their opinions, focusing on the period between the 1990--1992 and 2005--2009 surveys. During this period, voters increasingly self-identified based on their \textit{economic} views. This pattern contrasts with the post-2005 trend shown in Figure \ref{fig:heat_map}, where voters increasingly self-identified based on their cultural views. Second, Panel~B in Appendix Figure~\ref{fig:opinion_position_cor} extends the analysis in Panel~A of the same figure to earlier periods. We find that the correlation between voters' cultural opinions and the cultural positions of the parties they support remained relatively stable between 1990 and 2009 and began to rise significantly only in the most recent wave.\footnote{Throughout this analysis, we remove countries and variables where we do not have sufficient coverage in earlier waves. See Data Appendix \ref{subsec:data_processing} for details. We replicate the original analysis with this subset of countries and variables in Panel~B of both Appendix Figures \ref{fig:heatmap_1_4_subset} and \ref{fig:opinion_position_cor}, and find similar results.} 
\endQuotable{}

\beginQuotable[salience]
We next study whether the change in voter priorities could be due to parties emphasizing cultural issues. Appendix Figure~\ref{fig:change_salience} shows the share of the manifestos devoted to cultural and economic issues, aggregating both conservative and progressive positions. While the salience of both issues has moderately increased over time, economic issues remain far more salient in the manifestos. The emphasis on cultural issues relative to economic issues has not increased over time. 
\endQuotable{}

\paragraph{Discussion} 
\beginQuotable[timing_discussion] 
Finding that voters' prioritization of cultural issues occurred after 2009 is consistent with explanations that emphasize recent changes driving PRRP support, i.e., specific shocks making cultural cleavages more salient \citep{Bonomi2021}. Examples include changes in the media environment \citep{guriev3GInternetConfidence2020, tesei2022}, the financial crisis and austerity policies \citep{fetzerDidAusterityCause2019}, or the China shock \citep{gennaioli2025presidential}.
\endQuotable{} Finding that parties did not prioritize cultural issues over time compared to economic issues is less consistent with theories of political entrepreneurs making cultural issues prominent, to the extent that manifestos capture such changes. \beginQuotable[timing_caveat] However, we note that the analysis in this section is based only on the timing of different changes and should therefore be interpreted as suggestive.
\endQuotable{}

\vspace{0.5cm}
\subsection{Changes in Voter Attributes} \label{subsec:attributes}

Our finding that changes in voter attributes explain a relatively small share of the increase in PRRP support challenges common narratives of a swift shift in public opinion.  In this section, we explore this finding in more detail using reduced-form analyses of the IVS data. 

Figure \ref{fig:populist_score_agg} shows that, on average, voters did not move closer to the positions of PRRP since 2005. The top left panel shows the changes in the PRRP score defined in Section \ref{sec:data-attributes}, which aggregates opinions that are associated with PRRP support at the individual level. The thick blue line shows the average value across all countries, whereas the thin gray lines show the trends in each country. While there are some significant idiosyncratic changes in public opinion in specific countries, there has not been a substantial shift in voter attributes toward opinions associated with PRRP on aggregate.
The difference in average PRRP score between the 2005--2009 and 2017--2020 waves is only \lastWavePrrpScore{} standard deviations, which is \trendScoreShareOfDiffUnweighted{}  of the difference between the average PRRP score of PRRP supporters and other voters. 


Examining trends in specific voter opinions yields similar results. The remaining eight panels of Figure \ref{fig:populist_score_agg} plot the change in the eight opinion variables that are most strongly associated with the PRRP score, in each country and in aggregate.
%\footnote{The correlations are shown in Appendix Figure \ref{fig:populist_score_cor}.} 
Once again, we find that on average, voter opinions on these issues are relatively stable. For some of these issues, opinions have even moved away from the populist radical right.

Moreover, the rise in PRRP support has been documented not only in countries where opinions are shifting toward the populist radical right (e.g., Hungary) but also in countries where they are shifting away from it (e.g., Germany). Indeed, Appendix Figure \ref{fig:corr_score_support} shows that with the exception of Hungary, there is very little correlation between the change in the PRRP score and the change in PRRP support between 2005 and 2017. This is why our decomposition analysis finds that voter attributes do not drive the rise of PRRP. This finding demonstrates the importance of broad comparative perspectives in the study of electoral developments such as the rise of PRRP.

Hypothetically, even if voter opinions are stable on average, support for PRRP could still increase if opinions became more polarized, shifting to the right among some voters and to the left among others. Appendix Figure \ref{fig:populist_score_country} presents the distribution of PRRP scores for each country in both the 2005--2009 and 2017--2020 waves and does not find evidence of substantial polarization within countries. The PRRP score increases homogeneously in some countries (e.g., Hungary) and decreases homogeneously in others (e.g., Germany). 

% Can shorten in the future
Which voter opinions change? In Appendix Figure \ref{fig:opinions_change}, we take into account all IVS opinion variables and present the opinions that changed the most between the 2005--2009 and 2017--2020 survey waves. Interestingly, some of the most important changes move in the opposite direction of PRRP positions. For example, voters developed less conservative opinions on traditional morality issues, such as divorce and abortion. This result is consistent with our decomposition finding that the rise of PRRP is not driven by a shift in public opinion.

\subsubsection{The Reservoir of Potential PRRP Voters across Countries \label{subsec:cross_country}}
The stability of public opinion on issues associated with PRRP suggests that even before the electoral success of PRRP, there was a preexisting reservoir of potential PRRP voters. The size of this reservoir varies cross-nationally, which could explain why these parties are more successful in some countries than in others. However, measuring the number of potential PRRP voters across countries is challenging because voters in different countries face different sets of parties within different political systems. 

In this section, we use a similar decomposition exercise, across space instead of over time, to measure the size of the PRRP reservoir across countries. We focus on the 2017--2020 survey wave. For voters in each country, we simulate their counterfactual support for the National Front if they faced the same choice set as voters in France (the party was renamed the National Rally in 2018). In other words, we fix the set of parties, their characteristics, and their residuals to the values in France and allow only the voter attributes to differ across countries. Fixing all other components allows us to consistently compare the size of the reservoir of potential PRRP voters across countries. Formally, we calculate the following counterfactual for every country ($c$) in the 2017--2020 IVS wave based on Equation \ref{eq:prrp_share}: 

\begin{equation} \label{eq:counterfactual_france}
    \Tilde{S}^{c}_P = \int P(\Pi|x_{i};(\Phi,\beta),Z^{France},\zeta^{France})f^c(x_{i})\text{d}x_{i}.
\end{equation}

We find substantial cross-national variation in the size of the reservoir of potential PRRP voters. Figure \ref{fig:counterfactual_france} presents the counterfactual support for the National Front and reveals geographical patterns that are consistent with the actual support for PRRP shown in Figure \ref{fig:populist_support_map}. Potential support for PRRP is smallest in the Nordic countries and largest in Eastern Europe. While differences in voter attributes do not explain the increase in support for PRRP over time, they do explain some of the variation in support across countries. 

In some countries, the electoral system limits the number of parties, and therefore, the full potential support of PRRP may not be observed directly. Figure \ref{fig:counterfactual_france} shows that our method can calculate the size of the reservoir of potential PRRP supporters even in such countries. For example, we predict a much larger share of counterfactual support for PRRP in the U.K. than what is observed in practice. This is likely due to the first-past-the-post system, which creates higher entry barriers into parliament \citep{norris2005radical,fujiwara2011regression}. 

Finally, our method allows us to calculate the reservoir of potential PRRP voters in non-European countries, such as the United States.
Even though the two-party system in the U.S. impedes PRRP from gaining votes, we can estimate their latent support. We find that the PRRP reservoir is relatively large but smaller than in Eastern Europe. These results suggest that the U.K. and U.S. both have a large reservoir of potential PRRP voters within the voter base of other parties. 

Since voter attributes generate variation in PRRP support across countries, our finding that voter attributes do not drive the rise of PRRP is unlikely to stem from an omission of attributes that predict PRRP support. As a robustness test, Appendix Figure \ref{fig:counterfactual_germany} conducts a similar exercise using German parties and presents the counterfactual support for the AfD. The results are similar, demonstrating that our finding is not unique to France. 

\subsection{Changes in Party Characteristics}\label{subsec:party_positions}
Our decomposition estimates reveal that changes in party characteristics contribute rather little to the rise in PRRP support since 2005. 
In this section, we further investigate changes in party characteristics over time. Figure \ref{fig:prr_index_change} presents average changes in the two CMP indices for our five main party families: PRRP, center-left, center-right, Green, and socialists. Within each party family and year, we calculate the weighted average of each index. We weight each party by its vote share within its country and then weight all countries equally. Since elections are typically held every few years, we present five-year moving averages. 

We do not find that party positions changed in a way that is consistent with a rise in PRRP support. First, previous work has found that in the 1990s, the economic positions of mainstream parties converged, which could have pushed some dissatisfied voters toward PRRP \citep{berman2021causes}. 
%This may have pushed some dissatisfied left-wing voters toward PRRP \citep{berman2021causes}. 
However, the top panel of Figure \ref{fig:prr_index_change} shows that in the time period we study, the economic positions of the center-left and center-right parties have, in fact, slightly diverged.
%shows a divergence in economic positions, which are driven by the left and the greens. 
%We find that mainstream right parties and PRRP have similar economic positions and that, unsurprisingly, these positions are more conservative than green parties and left parties. 
Second, there is no evidence that PRRP gained votes by moderating their positions. On economic issues, PRRP positions remained stable over time. On cultural issues, which particularly distinguish PRRP from the other party families, PRRP became even more extreme, as shown in the bottom panel of Figure \ref{fig:prr_index_change}.\footnote{To better understand the shift of PRRP to the cultural right, Appendix Figure \ref{fig:prr_elm_change} shows the trends for the six party positions with the largest differences between PRRP and non-PRRP. The most substantial change occurred in positive mentions of a national way of life. PRRP today dedicate around 10\% of their manifestos to this issue, while the other parties dedicate approximately 1\%--3\%.} 
Third, we find some support for the claim that left-wing parties shifted further to the left on cultural issues, which may in turn have alienated working-class voters.  However, the substantial shift to the left occurred among socialist parties, and we do not see a similar shift among more centrist parties. Therefore, it is unlikely that this change would dramatically increase support for PRRP.\footnote{One concern with Figure \ref{fig:prr_index_change} is that the average of each party's position could also change as a result of a change in the vote share, which we use as weights within countries.  Therefore, Appendix Figure \ref{fig:position_change_fix_wgt} repeats the same exercise, allowing only the party positions to change over time. For every five-year period, we estimate the average change in each party's positions, weighting parties within countries by their initial support in that five-year period (while keeping equal weights for each country). We then plot the cumulative change for each party family. The results are similar to those reported in Figure \ref{fig:prr_index_change}, suggesting that the trends we find in party positions are not due to change in party support. } 



\beginQuotable[SupplyDemographics]

It is also possible that a change in PRRP leadership increased their support. For example, nominating more women may broaden support for these male-dominated parties \citep{weeks2023mannerparteien, AdamsBrackenGidronHorneOBrienSenk2023}. However, Figure \ref{fig:change_leader_demo} shows that the demographic composition of PRRP leadership has not changed dramatically over time. We present changes in three demographic indices by party family: the share of male legislators, the share of older (fifty years and above) white men among party leaders, and the share of non-establishment leaders, defined as those who are divorced, unmarried, do not hold a college degree, or had a below-median-wage occupation before entering politics. Across all three measures, PRRP values are fairly stable.

\endQuotable

To conclude, the reduced-form results are inconsistent with the major supply-side hypotheses. This provides intuition for why our decomposition suggests that, on aggregate, changes in party characteristics do not drive the increase in PRRP support.

\subsection{Changes in Residuals} \label{sec:residual}

The second major component in our decomposition is the residual, which accounts for \perResidualAllTime{} of the overall increase in support for PRRP. As discussed in Section \ref{subsec:Dynamic-Analysis-Decomposing}, this component captures factors such as changes in unobserved party characteristics, model misspecification, and the entry of new parties. 

Appendix \ref{sec:residual_appendix} examines the rise in the residual component in more detail. We show that most of the increase in this component is driven by party entry. \beginQuotable[EntryFigureText]We further decompose the residual component into the residuals of new and existing PRRP.\footnote{We calculate the same counterfactual as before, where only the residual component changes over time, and other components are held fixed at their initial levels in 2005--2009. However, we now separate the overall support for PRRP into support for new PRRP entrants and PRRP that already existed in our data in the 2005--2009 wave. See Appendix \ref{subsec:residual_new} for details.} Appendix Figure \ref{fig:dec_new_entry} shows that the rise of the residual component is more than fully driven by new entrants.\endQuotable{}

\beginQuotable[residualMain]
Therefore, we next analyze possible explanations for the substantial rise in support for new PRRP entrants. We argue that the entry of parties is unlikely to reflect a simple supply shock, in which voters were offered the option to vote for a PRRP for the first time. This supply-side explanation is unlikely since Table \ref{tab:party_entry_list} shows that in all countries with PRRP entry, either the same party or another radical right party already existed during the previous survey wave. We also present evidence that demand shocks are not the primary driver of new PRRP entry. Using counterfactual analyses, we show that substantial latent demand for PRRP entrants was already present prior to their observed entry.


Finally, we propose two explanations for the rise of PRRP entry that are consistent with the data: strategic considerations and unobserved support. First, the large increase in support for new PRRP entrants may reflect parties overcoming strategic considerations or coordination failures that prevented them from receiving substantial support in the past. For example, voters may coalesce around a party only when they anticipate that the party will receive sufficient support and pass the threshold needed to enter parliament \citep{freden2014threshold,freden2017opinion}.\footnote{Strategic considerations could interact with a change in voter preferences. It is possible that due to the change in voter priorities, the latent support for potential PRRP entrants passed some threshold and, as a result, strategic considerations led to a substantial increase in the parties' observed support and they entered the political system. This would mean that some of the residual component reflects a change in priorities.}
Second, we find that at least some of the support for new PRRP entrants already existed in previous waves but did not appear in our data. This is either because support for small parties is not always recorded in the IVS or because of mergers and splits with parties not classified as PRRP. 
While unobserved prior PRRP support increases the residual component, as entry is always attributed to the residual, this does not affect our main conclusion. Since all the additional growth is attributed to the residual, unobserved prior support does not change the \emph{relative} importance of our three main components (voter priorities, voter attributes, and party characteristics).

Strategic considerations and unobserved support could similarly affect non-PRRP entrants. Therefore, we examine whether cases where parties have latent demand but still do not enter exist for non-PRRP as well. We use our model to calculate the latent demand for all new parties, both PRRP and non-PRRP, before they appear in our data. We find similar pre-entry latent demand across entrants of both types. The similar latent demand for PRRP and non-PRRP entrants is more consistent with explanations common to all new parties (strategic considerations, unobserved support) than with explanations unique to PRRP, such as their normalization. 

\endQuotable{}

\section{Conclusions}\label{sec:conclusion}

There is no lack of explanations for the rise of the populist radical right. Our goal in this manuscript is not to introduce another factor that may have contributed to these parties' electoral success, but rather to assess the explanatory power of leading categories of existing explanations. We find that changes in priorities provide the strongest explanation for the growing electoral appeal of PRRP. We present comprehensive empirical evidence supporting Bartels's \citeyearpar{bartels2023democracy} memorable argument that PRRP are not surfing into power on a wave of growing nativism and authoritarianism in public opinion; instead, these parties have proved adept at mobilizing preexisting reservoirs of potential support. This implies that significant electoral changes can occur not only when people change their minds but also when certain issues become more consequential in shaping vote choice.
Understanding the sources of this shift in voter priorities is a promising path for future research.

%\newpage
%\clearpage
\singlespacing



% Remove word "in" in bibliography 
\renewbibmacro{in:}{}

% remove pp
\DeclareFieldFormat{pages}{#1}

%\end{refsection}
\printbibliography

%====================================== FIGURES AND TABLES =======================================%

\clearpage
\section*{Tables and Figures}


\newcommand{\figureMacro}[4] {
    \begin{figure}[!htp]
    \caption{\normalsize #1}\label{fig:#2}
    \includegraphics{output/figures/#3} 
    \floatfoot{#4}
    \end{figure}
    %\clearpage
}



\newcommand{\subfigureMacro}[3] {
    \begin{subfigure}[!htp]{1\textwidth}
    \centering
    \caption{\normalsize #1}\label{fig:#2}
    \includegraphics{output/figures/#3} 
    \end{subfigure}
}


\newcommand{\tableMacro}[4] {
    \begin{table}[h]
    \caption{#1}\label{tab:#2}
    \input{output/tables/#3} 
    \vspace{-3mm}
    \floatfoot{#4}
    \end{table}
 %   \clearpage
}




%====================FIGURES ======================%

\begin{figure}[!htp]
    \caption{\normalsize Support for Populist Radical Right Parties Over Time}\label{fig:populist_support}
    \includegraphics[width=0.65\textwidth]{output/figures/manifesto_populist_rise/fr_vote_share_global.pdf} 
    \floatfoot{\textit{Notes:} This figure shows the average vote share of PRRP in the  \countriesNum{} European countries in our data. Within each country, the vote share every year is calculated as the average PRRP vote share among all parties appearing in the CMP dataset in all parliamentary elections in the five-year window centered around that year. The figure shows the average share across all \countriesNum{} countries. The gray bars mark when the surveys of the three IVS waves used throughout the paper were held.}
\end{figure}

\begin{figure}[!htp]
    \caption{\normalsize Support for Populist Radical Right Parties by Country, 2017--2020 IVS Wave}\label{fig:populist_support_map}
    \includegraphics[width=0.9\textwidth]{output/figures/WVS/populist_support_maps/prrp_support_map_wave6.pdf} 
    \vspace{-0.5cm}
    \floatfoot{\textit{Notes:} This figure shows the average PRRP support in the 2017--2020 IVS wave for the \countriesNum{} countries covered in our analysis. PRRP support is calculated as a share of support for all parties appearing in the data. IVS sampling weights are used.}
\end{figure}


\begin{figure}[!htp]
    \caption{\normalsize Voting Weights Distribution by Supported Party Family, 2017--2020}\label{fig:weights_static}
    \includegraphics{output/figures/model_weight/weight_static_index_and_highbtwgrvar_weighted_sd.pdf} 
    \vspace{-0.5cm}
    \floatfoot{\textit{Notes:} This figure shows the distribution of weights voters place on the economic and cultural position indices and eight individual manifesto positions in the most recent survey wave (2017--2020) separately for supporters of different party families. For each combination of party family and index/position, we present the 25th, 50th, and 75th percentiles of the corresponding weight distribution. We estimate the model using the 2017--2020 survey wave and calculate each voter's weights based on her attributes using Equation~\ref{eq:weights}. Weights are measured in units of standard deviation to utility units, i.e., the increase in utility for a one standard deviation increase in the position. We aggregate individual weights into indices based on the CMP economic and cultural indices, described in Section \ref{subsec:CMP}. Specifically, we calculate the change in the utility for a one standard deviation increase in the index, generated by equal-sized changes in all positions used in the CMP index; weights on left-wing positions are multiplied by -1. The eight individual positions presented are the ones with the largest variance in weight between supporters of the different party families. \ivsWeights{} The party families are described in Appendix \ref{subsubsec:classification}.}
\end{figure}

\clearpage



\figureMacro{Decomposition of the Rise in Populist Radical Right Support}
{decomposition}
{decomposition/waterfall_total_46.pdf}
{\textit{Notes:} This figure presents the results of our main decomposition exercise. The black bars mark the overall increase in PRRP support between 2005--2009 and 2017--2020, which we set to 100\%. We aggregate across all countries with PRRP support in our data using a weighted average of each country's decomposition results (see Appendix Table \ref{tab:decomposition_table}). 
Weights are the inverse of the share of PRRP support in the 2017--2020 wave. We use IVS sampling weights within countries.}

\beginQuotable[fig_dec_robust]
% Tag the figure block so other TeX files can import it directly.
%<*figureTest>
\begin{figure}[!htp]
    % Keep the original figure content unchanged for the main paper.
    \caption{\normalsize Decomposition Robustness}\label{fig:decomposition_robustness}
    \includegraphics{output/figures/decomposition/Decomposition_Robustness_Waterfalls.pdf} 
    \floatfoot{\textit{Notes:} These figures plot the decomposition results for changes in PRRP support between the 2005--2009 and 2017--2020 survey waves under different specifications. Panel~A presents the Shorrocks Index, an average of different decomposition orders. Panel~B presents results from a leave-country-out exercise, where for each country, we estimate the model parameters using data from all other countries. Panel C shows the decomposition results when we use a bliss-point model, as described in Appendix \ref{sec:bliss}. Panel D shows the results when we replace voters' ages with their year of birth (cohort). Panel E repeats the entire exercise with additional supply-side variables on parties' rhetoric (complexity and insult score). Panel F repeats the exercise with party positions based on pledges instead of the full manifestos. Panels G and H show results from estimating voter priorities and running the decomposition separately for higher-income and lower-income countries. The higher-income countries are Austria, Denmark, Finland, France, Germany, Iceland, the Netherlands, Norway, Sweden, Switzerland, and the UK. The lower-income countries are Estonia, Greece, Hungary, Italy, Lithuania, Poland, Portugal, Slovenia, and Spain. All figures aggregate results across countries, weighting each country by the inverse share of PRRP support in the 2017--2020 wave. We use the IVS sampling weights within countries. See further details in Section \ref{sec:decomp_robustness}.}
    \end{figure}
%</figureTest>
\endQuotable{}
\clearpage



\figureMacro{Changes in the Distribution of Voting Weights, Holding Attributes Fixed}
{weights_dynamic}
{model_weight/weight_distr_index_allpar_standardized_econ_soc_46.pdf}
{\textit{Notes:} This figure shows the distribution of the weights voters place on the economic and cultural indices of party positions for different waves. We use the CMP economic and cultural indices described in Section \ref{subsec:CMP}; their manifesto components are listed in Appendix Table \ref{tab:manifesto_vars}. Weights are calculated based on Equation \ref{eq:weights} using the estimated voter priorities (utility parameters) for the 2005--2009 wave (solid blue lines) and the 2017--2020 wave (dashed yellow lines). The voter attributes are fixed at their distribution in 2017--2020, so that weights only change due to a change in priorities. The weight placed on an index is the change in the utility for an increase in one standard deviation in the index, driven by an equal-sized change in each party position included in the CMP index, where weights for positions that enter the index with a negative sign are multiplied by -1. \ivsWeights{}}




\figureMacro{Change in  Predicted Ideological Self-Identification}
{heat_map}
{heat_map_voters/voters_heatmap_4to6_rescale.pdf}
{\textit{Notes:} The figure shows the change in voters' predicted ideological self-identification, on a 1--10 scale, between the 2005--2009 survey wave and the 2017--2020 survey wave as a function of voters' economic and cultural opinions. Economic and cultural opinions are aggregated into two indices using the first principal component. The PCA loadings for the economic and cultural opinion indices are estimated once using the pooled 2005--2009 and 2017--2020 sample. The list of IVS variables included in each index appears in Appendix Table~\ref{tab:ivs_vars}. We use a nonparametric regression model to predict ideological self-identification (higher values are associated with more right-wing ideology), based on the voter's economic and cultural opinion indices, for each survey wave separately. We then subtract the prediction for the earlier wave from that for the later wave. The two contour lines mark the values of  -0.1 and 0.1.
}

\beginQuotable[figPriorHet]
\begin{figure}[!htp]
    \caption{\normalsize Relative Increase in Support for PRRP due to Change in Priorities}\label{fig:hetero_prior_prrp}
    \includegraphics{output/figures/decomposition/weighted_avg_diff_from_full_pop_bar.pdf} 
    \floatfoot{\textit{Notes:} This figure presents the counterfactual change in support for PRRP when only voter priorities change, for different subsets of the population. We calculate the counterfactual rise in support for PRRP when only priorities change, while all other components are held fixed at their values in 2017--2020 (the first line of Equation \ref{eq:decomposition}) for various subgroups. \ivsWeightsCountryInverse{} We then compare the rise in support for PRRP in every subgroup relative to the rise in support among the full population. \PRRPscore{} 
    }
\end{figure}
\endQuotable{}





\clearpage



\figureMacro{Evolution of Voter Opinions}
{populist_score_agg}
{WVS/populist_voters/PRRP_pop_opinion_over_time.pdf}
{\textit{Notes:} This figure presents the PRRP score and the average values of the eight opinion variables most strongly associated with PRRP support, by survey wave. The thin gray lines show the trends in each country, while the thick blue lines show the average across all \countriesNum{} countries.  \PRRPscore Among the IVS opinion variables, we select the eight variables with the largest absolute correlation with the PRRP score. We standardize the eight opinion variables within each country using standard deviations from the 2005--2009 IVS wave. For any opinion that is unavailable in a country in 2005--2009, that country is omitted from the corresponding panel. In all panels, we set the mean value in the 2005--2009 wave to zero in all countries. For more details on each variable, see Appendix Table \ref{tab:ivs_vars}. \ivsWeights{}}





\figureMacro{Counterfactual Support for the National Front by Voter Attributes}
{counterfactual_france}
{counterfactual_maps/france_counterfactual_with_US.pdf}
{\textit{Notes:} This figure shows the counterfactual support for the National Front in the 2017--2020 wave if French voters had the attributes of voters in other countries. We calculate the counterfactual support for the National Front separately for each country based on the formula in Equation \ref{eq:counterfactual_france}. In all countries, we use the party characteristics of French parties in the 2017--2020 wave ($Z^{France}$) along with the estimated residuals for French parties  ($\zeta^{France}$) and the voter priorities that were estimated for this wave ($\hat{\Phi},\hat{\beta}$). For each country, we predict the share of National Front supporters according to the voter attributes in that country ($f^c(x_i)$). For the U.S. sample, we impute the responses to one IVS question about the European Union based on the sample average. For all countries, we use IVS sampling weights.}



\figureMacro{Changes in Party Positions Over Time}
{prr_index_change}
{manifesto_party_pos_change/index_aggregate.pdf}
{\textit{Notes:} This figure shows the changes in the CMP economic and cultural indices for five party families--PRRP, center-right, center-left, socialist, and green. Indices measure the net share of the manifesto dedicated to right-wing positions. Each index is constructed by adding the manifesto shares of conservative positions and subtracting the shares of liberal positions such that positive values reflect more support for a free market or more conservative cultural values. The figure presents five-year moving averages for each index and each party family. Each country is weighted equally, and parties within each country are weighted by their vote share. The indices are discussed in Section \ref{subsec:CMP}, their manifesto components are described in Appendix Table \ref{tab:manifesto_vars}, and the party classification is described in Appendix \ref{subsubsec:classification}.} 


\clearpage


%====================TABLES ======================%



\tableMacro{Descriptive Statistics}{sum_stat}
{top5_manifesto_wvs.tex}
{\textit{Notes:} This table provides descriptive statistics on party positions and leadership demographics (Panel~A) and voter attributes based on IVS data (Panel~B). The first two columns present the average of each variable in the 2005--2009 wave and the last two columns present their average in the 2017--2020 wave. 
In Panel~A, the first two rows show the averages of the CMP economic and cultural indices.  The indices measure net shares, i.e., the difference between the share of platforms dedicated to right-wing positions and left-wing positions. See Appendix Table \ref{tab:manifesto_vars} for further details on each position, as well as a list of variables included in the two indices. The next five rows show the positions with the largest difference between PRRP and non-PRRP across both waves. 
%To pick them, we standardize all positions with parties weighted by vote share within each country and all countries weighted equally. We then calculate the average value for PRRP and non-PRRP in each position giving each wave equal weight. We pick the five positions with the largest standardized difference in absolute terms.
Each variable represents the share of the platform mentioning that position in percentage terms. The last three rows represent demographic statistics for party leadership. In Panel~B, the first five rows show key demographic variables.  We define urban as living in a city with more than 100,000 people. The next five rows focus on the opinions with the largest difference between PRRP and non-PRRP supporters. Opinions are measured in units of standard deviations from the mean. In both panels, averages are calculated with equal weight for each country. In Panel~A, within countries, parties are weighted based on their vote share. In Panel~B, within countries, we use the IVS sampling weights.}
\clearpage




%====================================== APPENDIX =======================================%

\clearpage
\appendix
\renewcommand{\thefigure}{A.\arabic{figure}}
\setcounter{figure}{0}     
\renewcommand{\thetable}{A.\arabic{table}}
\setcounter{table}{0}  

\part*{Appendix}
\section{Additional Figures and Tables}

\begin{refsection}





%====================================== APPENDIX FIGURES =======================================%
\beginQuotable[flow_chart]

\begin{figure}[!htp]
\centering
\input{code/documents/input_common/model_flow}

\caption{Voting Model Summarized}
\label{fig:model_flow}

\floatfoot{\textit{Notes:} This diagram summarizes the model in Section \ref{sec:model}. We show how each of the four inputs, voter attributes, voter priorities, party characteristics, and the residuals, are combined to determine voter utility, and specify the source of each input. }

\end{figure}

\endQuotable{}

\clearpage


\figureMacro{Covariates Most Strongly Correlated with the PRRP Score}
{populist_score_cor}
{WVS/populist_voters/PRRP_cor_weighted.pdf}
{\textit{Notes:} This figure presents the voter attributes most strongly correlated with the PRRP score. The score is calculated by running a LASSO regression predicting PRRP support. The regression is run on the 2017--2020 survey wave and includes all IVS variables in our final dataset, along with country fixed effects, with no penalty on the country fixed effects. To calculate the PRRP score, we standardize the fitted value generated by the regression coefficients (excluding country fixed effects) with each country weighted equally. \ivsWeights{}}
\endQuotable{}
\clearpage

\begin{figure}[!htp]
\thispagestyle{empty}
\includegraphics[width=1.12\textwidth, angle =90 ]{output/figures/WVS/voter_correlations/correlations_average_diff.pdf} 
\caption{\normalsize Similarities between Party Voters \label{fig:party_correlation}}
\vspace{-0.5cm}
\floatfoot{\textit{Notes:} This figure presents the similarities between voters of different parties across countries. Red colors indicate a positive correlation and blue colors indicate a negative correlation. For each party, we use a linear regression to predict support for that party based on voter attributes, using data from the 2017--2020 IVS wave in the relevant country. Then, for each pair of parties from different countries, we calculate the correlation between the fitted values across all voters in the two countries. For example, to find the correlation between AfD and the National Front, we estimate the correlation between the fitted probability that voters in both Germany and France would support AfD and the fitted probability that those voters would support the National Front, where both fitted probabilities are linear predictions based only on their attributes. The labels show the average correlation between all parties in each family. We determine whether a party is PRRP based on the PopuList dataset classification of radical-right parties. We use the PopuList definition of populism for the ``other populist parties'' category. We classify the remaining parties into party families based on the CMP data as explained in Section \ref{subsec:party_class}. The figure is based on all parties in our data that received support from over 50 respondents. IVS sampling weights are used in all regressions. More details appear in Appendix \ref{sec:classify_right}.}
\end{figure}
\clearpage

\beginQuotable[fig_phi_all]
\begin{figure}[p]
\caption{Voter Attributes With the Largest Coefficients in the Weight Indices \label{fig:phi_all}}
\begin{subfigure}{1\textwidth} 
\centering
\caption{Economic and Cultural Indices \label{fig:phi_all_a}}
\label{fig:phi_pos}
\includegraphics[width=0.83\textwidth]{output/figures/phis/phi_newZs_pos.pdf}

\end{subfigure}

\vspace{5mm}

\begin{subfigure}{1\textwidth}
\centering
\caption{Leader Demographics Indices}
\label{fig:phi_demo}
\includegraphics[width=0.83\textwidth]{output/figures/phis/phi_newZs_demo.pdf}

\end{subfigure}


\floatfoot{\textit{Notes:} This figure shows the largest coefficients (in absolute terms) on the (standardized) IVS variables generating the weight indices. We calculate weights on individual party characteristics using Equation \ref{eq:weights}. Panel (a) aggregates individual weights into indices based on the economic and cultural party position indices suggested by the CMP (Section \ref{subsec:CMP}). Specifically, we take a simple average of the weights on all positions used in the CMP index, multiplied by -1 for left-wing positions. Panel (b) aggregates individual weights into indices based on leadership demographics (Section \ref{subsec:leadership}). The Older White Men Index is constructed as the sum of the weights on male and age minus the weight on minority; the Non-Establishment Index is constructed as the weight on the number of divorces minus the weights on years of schooling, married, and occupational wage. For each index, we present the ten largest coefficients in absolute value.}
\end{figure}

\endQuotable{}
\clearpage

\beginQuotable[decomposition_robustness_app]
\begin{figure}[!htp]
    \caption{\normalsize Additional Decomposition Robustness}\label{fig:decomposition_robustness_app}
    \includegraphics{output/figures/decomposition/Decomposition_Robustness_Waterfalls_Appendix.pdf} 
    \floatfoot{\textit{Notes:} This figure plots the decomposition results for changes in PRRP support between the 2005--2009 and 2017--2020 survey waves under different specifications. Panel~A shows our main decomposition. Panel~B repeats the entire exercise, excluding voters' demographic variables and party leadership demographics. Panels C and D estimate the $\beta_t$ parameters (Equation \ref{eq:oaxaca}) using a lower ($k=3$) and higher ($k=7$) number of dimensions for $\beta_t$, compared to our original specification ($k=5$). Panels E and F repeat the exercise using different sets of countries, where Panel E adds the Czech Republic and Slovakia, and Panel F excludes Greece and Italy. All figures aggregate results across countries, weighting each country by the inverse share of PRRP support in the 2017--2020 wave. We use the IVS sampling weights within countries. See further details in Section \ref{sec:decomp_robustness}.}
    \end{figure}
    \endQuotable{}
\clearpage

\begin{figure}[ht]
\captionsetup{justification=centering}
\caption{ Changes in the Distribution of Voting Weights, Holding Attributes Fixed: \\ Alternative Economic and Cultural Indices \label{fig:alternative_cmp}}
\subfigureMacro{Bakker and Hobolt (2013)}
{bakker_hobolt}
{model_weight/alternative_cmp_measures/bak_hob_weight_distr.pdf}

\vspace{5mm}
\medskip
\subfigureMacro{Prosser (2014)}
{prosser}
{model_weight/alternative_cmp_measures/prosser_weight_distr.pdf}

\captionsetup{justification=justified}

\floatfoot{\textit{Notes:} This figure shows the distributions of the weights voters place on two alternative economic and cultural party position indices proposed by (a) \citet{bakker2013measuring} and (b) \citet{prosser2014building}. Weights are calculated based on Equation \ref{eq:weights} using the estimated voter priorities for the 2005--2009 wave (solid blue lines) and 2017--2020 wave (dashed yellow lines). The voter attributes are fixed at their distribution in 2017--2020 so that weights only change due to a change in priorities. The weight placed on an index is the average weight corresponding to each party position included in the index, where weights for positions that enter the index with a negative sign are multiplied by -1. Appendix Table \ref{tab:manifesto_vars} describes the variables included in each index. Weights are measured in units of standard deviation to utility, i.e., the increase in utility for an increase of one standard deviation in the index. \ivsWeights{}}
\end{figure}
\clearpage

\beginQuotable[fig_altInd]
\begin{figure}[!htp]
\captionsetup{justification=centering}
\caption{Changes in the Distribution of Voting Weights, Holding Attributes Fixed: \\Indices on Additional Topics \label{fig:alternative_indices}}
\includegraphics{output/figures/model_weight/weight_distr_index_allpar_standardized_ex_index_46.pdf}
\captionsetup{justification=justified}
\floatfoot{{\textit{Notes:} This figure shows the distribution of the weights voters place on the expanded hawk-dove \citep[based on][]{kiratli2021hawkdove}, European integration, and government efficiency indices of party positions for different waves, holding the distribution of voter attributes fixed. Positive values represent support for hawkish positions, EU disintegration, and greater government efficiency, and negative values represent support for the opposing positions. The manifesto components of these alternative CMP indices are described in Appendix Table \ref{tab:manifesto_vars}. Weights are calculated based on Equation \ref{eq:weights} using the estimated voter priorities for the 2005--2009 wave (solid blue lines) and 2017--2020 wave (dashed yellow lines). The voter attributes are fixed at their distribution in 2017--2020 so that weights only change due to a change in priorities. The weight placed on an index is the average weight corresponding to each party position included in the CMP index, where weights for positions that enter the index with a negative sign are multiplied by -1. Weights are measured in units of standard deviation to utility units, i.e., the increase in utility for an increase of one standard deviation in the index. \ivsWeights{}}}
\end{figure}
\endQuotable{}
\clearpage


\beginQuotable[fig_Germany]
\figureMacro{Case Study: Most Important Issues in Germany Over Time}
{germany_case_study}
{case_studies/Germany/stacked_bar_plot_categorized.pdf}
{\textit{Notes:} This figure shows the distribution of answers to the ``most important issue'' question in the German Politbarometer, translated as: ``In your opinion, what is currently the most important problem in Germany?'' We include the following issues in the Economics category: ``Economic situation, growth, recession'', ``Unemployment, jobs'', ``Prices, inflation, cost of living'', ``Wages and income, wage negotiations'', ``Taxes, tax reform, fees'', ``Pension provision, pension issues'', ``social divide/ justice/ rich-poor'', ``National debt'', ``International banking crisis''. The following issues are included in the Culture (ex. immigration) category: ``family / children / youth'', ``Violence, extremism (left/right), demonstrations, Antifa'', ``School \& Education'', ``Environment, climate change, renewable energy''. The issue ``Immigration, integration, asylum, refugees'' was categorized separately as ``Immigration''. The rest of the issues were categorized as ``Other'': ``Foreign policy, general'', ``Federal elections outcome'', ``Healthcare/Health policy/Long-term care insurance'', ``Party/political dissatisfaction/politician affairs'', ``Viruses \& Pandemics'', ``Other problems'', ``Other''.}
\endQuotable{}
\clearpage

\beginQuotable[fig_corOld]

\begin{figure}[p]
    \captionsetup{justification=centering}
    \caption{Correlations of Voter Opinions and the Positions of their Preferred Party\label{fig:opinion_position_cor}}
{\centering
\scalebox{0.95}{%  
    \subfigureMacro{Main survey waves, all data}
    {opinion_position_cor_main}  {voter_party_correlation/voterPartyCorPCA.pdf}
   } 
    \vspace{2mm} % Space between subfigures

\scalebox{0.95}{%
    \subfigureMacro{Earlier survey waves, subset of data}
    {opinion_position_cor_1346}
    {voter_party_correlation/voterPartyCorPCA_1346_combined.pdf}
    }
    \captionsetup{justification=justified} % Reset justification for floatfoot
   
    \floatfoot{\textit{Notes:} 
%No changes to the first part:
This figure displays binned scatter plots of the relationship between the opinions of voters and the positions of the parties they supported over time. \ivsWeights{} The positions of parties are based on the Comparative Manifesto Project economic and cultural indices. The components of these indices are shown in Appendix Table \ref{tab:manifesto_vars}. The voter opinion indices are created using a principal component analysis of IVS variables related to economics and culture. The variables included in the principal component analysis appear in Appendix Table \ref{tab:ivs_vars}. Panel (a) shows the change between 2005--2009 and 2017--2020. Panel (b) additionally includes the 1990--1992 and 1999--2000 survey waves. To ensure consistency in the set of countries and variables used across all waves, Panel (b) excludes the following variable from the economic index: \textit{Personal over government responsibility should be increased}; the following variables from the cultural index: \textit{Prostitution justifiable}, \textit{Like idea of army rule}, \textit{Like idea of democracy}; and \textit{Country needs strong leader}; and the following countries: Estonia, Lithuania, Switzerland, Poland, and Norway.
}}
\end{figure}

\endQuotable{}
\clearpage


\beginQuotable[Fig_PriorHetTotal]
\begin{figure}[!htp]
    \caption{\normalsize Change in Cultural and Economic Weight Ratio, Holding Attributes Fixed}\label{fig:hetero_prior_change}
    \includegraphics{output/figures/model_weight/dist_by_group/hetero_prior_change.pdf} 
    \floatfoot{\textit{Notes:} This figure presents the average change in the log of the ratio of absolute values of the cultural and economic weight indices (Equation \ref{eq:chane_ratio}) from the 2005--2009 wave to 2017--2020 wave by subgroups.  A positive value means that cultural positions became more important compared to economic positions.  We use the CMP economic and cultural indices that are described in Section \ref{subsec:CMP}. Weights are calculated based on Equation \ref{eq:weights} using the estimated voter priorities for the 2005--2009 wave and 2017--2020 wave. The voter attributes are fixed at their distribution in 2017--2020, so that weights only change due to a change in priorities. The weight placed on an index is the change in the utility for an increase in one standard deviation in the index, driven by an equal-sized change in each party position included in the CMP index, where weights for positions that enter the index with a negative sign are multiplied by -1. The first row shows the mean change for the full population. The remaining rows show the mean for each subgroup. Confidence intervals are at the 95\% level. We report on the right the p-value for the hypothesis of no difference between the two subgroups.  \PRRPscore \ivsWeights{}}
\end{figure}
\endQuotable{}



\beginQuotable[fig_PRRPDem]
\begin{figure}[!htp]
    \caption{\normalsize Average PRRP Score by Subpopulation}\label{fig:demo_prrp_score_bar}
    \includegraphics{output/figures/decomposition/demo_prrp_score_bar.pdf} 
    \floatfoot{\textit{Notes:} This figure reports the average PRRP score for each subgroup in the 2017--2020 wave. \PRRPscore  Bars show weighted means of this standardized PRRP score within each subgroup.  \ivsWeights{}}
\end{figure}
\endQuotable{}



\clearpage




%%%%%%%%%%%% EARLIER HEATMAPS %%%%%%%%%%%%%%%%%
\beginQuotable[fig_heatOld]
\begin{figure}[p]
\captionsetup{justification=centering}
\caption{Change in Predicted Ideological Self-Identification, Earlier Time Period \label{fig:heatmap_1_4_subset}}


\begin{subfigure}[!htp]{0.67\textwidth}
\centering
\caption{\normalsize Change between 1990--1992 and 2005--2009}\label{fig:heatmap_1_4_subset_only}
\includegraphics[width=\textwidth]{output/figures/heat_map_voters/voters_heatmap_1to4.pdf} 
\end{subfigure}


\vspace{2mm}
\medskip

\begin{subfigure}[!htp]{0.67\textwidth}
\centering
\caption{\normalsize Change between 2005--2009 and 2017--2020}\label{fig:heatmap_4_6_subset}
\includegraphics[width=\textwidth]{output/figures/heat_map_voters/voters_heatmap_46_V2.pdf} 
\end{subfigure}

\captionsetup{justification=justified}

\floatfoot{\textit{Notes:} The figure shows the change in the predicted ideological self-identification, on a 1--10 scale. Economic and cultural opinions are aggregated into two indices using the first principal component. We use a nonparametric regression model to predict ideological self-identification (higher values are associated with more right-wing ideology), based on the voter's economic and cultural opinion indices, for each survey wave separately. We then subtract the prediction of the earlier wave from the latter. To ensure consistency in the set of countries and variables used across all waves, we exclude the following variable from the economic index: \textit{Personal over government responsibility should be increased}; the following variables from the cultural index: \textit{Prostitution justifiable}, \textit{Like idea of army rule}, \textit{Like idea of democracy}, and \textit{Country needs strong leader}; and the following countries: Estonia, Lithuania, Switzerland, Poland, and Norway. Subfigure (a) shows the change between 1990--1992 and 2005--2009. Subfigure (b) shows the change between 2005--2009 and 2017--2020. This subfigure is similar to Figure \ref{fig:heat_map} but uses a subset of countries and variables.}

\end{figure}
\endQuotable{}





\beginQuotable[fig_partySalience]
\figureMacro{Changes in the Salience of Topics in Manifestos Over Time}
{change_salience}
{manifesto_party_pos_change/salience_grouped_broad.pdf}
{\textit{Notes:} The figure shows the changes in the share of the manifesto covering economic and cultural topics across all countries and parties since 1990. 
The figure presents five-year moving averages for each topic. Each country is weighted equally, and parties within each country are weighted by their vote share. The manifesto components are described in Appendix Table \ref{tab:manifesto_vars}.}
\endQuotable{}
\clearpage


\figureMacro{Changes in PRRP Score and PRRP Support, 2005--2017}
{corr_score_support}
{WVS/populist_voters/PRRP_score_corr_support_weighted.pdf}
{\textit{Notes:} This figure plots the changes in the voters' average  PRRP score and the changes in the PRRP voting share by country. Changes in PRRP score are between the 2005--2009 survey wave and the 2017--2020 survey wave. Changes in PRRP voting share are between 2005 and 2017. For both years (2005 and 2017), we take the average of all elections held in a five-year window centered around that year. \PRRPscore The change in PRRP score is normalized by the standard deviation of the score in the 2005--2009 survey for each country. IVS sampling weights are used. PRRP vote shares are calculated relative to all parties in the CMP dataset in parliamentary elections held within the five-year window. }
\clearpage


\figureMacro{Voters' PRRP Score by Country}
{populist_score_country}
{WVS/populist_voters/PRRP_score_dist_density_country_ceil_weighted.pdf}
{\textit{Notes:} This figure presents the distribution of voters' PRRP score by country and survey wave. \PRRPscore The vertical dashed lines mark the average values (normalized to 0) and the vertical dotted lines mark the minimum and maximum possible values. IVS sampling weights are used.}
\clearpage


\figureMacro{Opinions that Changed the Most, 2005--2020}
{opinions_change}
{WVS/opinions_trend/opinions_average_trend_1st_half.pdf}
{\textit{Notes:} This figure shows the opinions that changed the most between the 2005--2009  and 2017--2020 survey waves. We run a separate regression of every opinion variable on the survey year and country fixed effects. Opinions are measured in standard deviation units. Each dot represents the time coefficient in this regression. 95\% confidence intervals are reported. \ivsWeights{}}
\clearpage




\figureMacro{Counterfactual Support for the AfD by Voter Attributes}
{counterfactual_germany}
{counterfactual_maps/germ_counterfactual_with_US.pdf}
{\textit{Notes:} This figure shows the counterfactual support for the AfD in the 2017--2020 wave if German voters had the attributes of voters in other countries. We calculate the counterfactual support for the AfD separately for each country based on a formula similar to Equation \ref{eq:counterfactual_france}. We use IVS sampling weights. For more details, see Figure \ref{fig:counterfactual_france}.}
%In all countries, we use the party positions  ($Z^{Germany}_t$) of German parties in the 2017--2020 wave along with the estimated residuals for German parties  ($\zeta^{Germany}_t$) and the model parameters that were estimated for this wave ($\hat{\Phi}_t,\hat{\beta}_t$). For each country,  we predict the share of AfD supporters according to the voter attributes in that country. For the U.S. sample, we impute the responses for questions related to European topics (e.g., opinion on the European Union) based on the sample averages.}
\clearpage

%zzz this should be the same variables as Tab1 (maybe also fig3) and we should have a 6th variables. 
\figureMacro{Changes in Most Distinctive Party Positions Over Time}
{prr_elm_change}
{manifesto_party_pos_change/top_distinct_aggregate.pdf}
{\textit{Notes:} This figure shows the average position by party family for the six positions with the largest differences between PRRP and the other parties. Each position is measured as the average percentage of the platforms devoted to that position by parties in that family. The manifesto variables are described in Appendix Table \ref{tab:manifesto_vars}. The figure presents five-year moving averages for each position and each party family. Each country is weighted equally, and parties within each country are weighted by their vote shares.}
\clearpage

\figureMacro{Within-Party Position Changes}
{position_change_fix_wgt}
{manifesto_party_pos_change/index_parfam_within.pdf}
{\textit{Notes:} This figure shows the cumulative within-party changes in the CMP economic and cultural indices by party family. The manifesto components included in each index are described in Appendix Table \ref{tab:manifesto_vars}. The figure presents cumulative average changes within each party since 2005. In particular, for each five-year period, we first compute position changes at the party level. We then aggregate the change across parties and countries. Each country is weighted equally, and parties within each country are weighted by their average vote share in the initial period. Thus, changes in the vote share do not affect the change in positions between consecutive periods. We present the cumulative change aggregating all changes since 2005. Note that the set of parties is not identical across periods.}
\clearpage





\beginQuotable[fig_changeLeaderDem]
\begin{figure}[!htp]
    \centering
    \caption{\normalsize Changes in Leadership Demographics Over Time}
    \label{fig:change_leader_demo}
    \includegraphics[width=0.9\linewidth]{output/figures/femaleShare/demo_weighted_all.pdf}
    \floatfoot{\textit{Notes:} This figure shows the changes in the share of male legislators, the probability of having an older white male leader, and the probability of having a non-establishment leader for five party families--PRRP, center-right, center-left, socialist, and green--since 2005. An older white male leader is defined as a leader over fifty years old, male, and from a non-minority group. A non-establishment leader is defined as one who is divorced or unmarried, does not hold at least a college degree, or worked in an occupation with below-median weekly wages before entering politics. The figure presents five-year moving averages for each index and each party family. Each country is weighted equally, and parties within each country are weighted by their vote share. Details on leadership demographics are provided in Table \ref{tab:leader_code_book}.} 
\end{figure}
\endQuotable{} 
\clearpage

%\figureMacro{Changes in Leader Demographics over Time}
%{change_leader_demo}
%{femaleShare/demo_weighted_all.pdf}
%{\textit{Notes:} This figure shows the changes in the share of male legislators, the probability of having an older white male leader, and the probability of having a non-establishment leader for five party families--PRRP, center-right, center-left, socialist, and green--since 2005. An older white male leader is defined as a leader over fifty years old, male, and from a non-minority group. A non-establishment leader is defined as one who does not belong to the group of leaders with at least a bachelor's degree, are married, not divorced, and held an occupation with above median weekly wages before entering politics. The figure presents five-year moving averages for each index and each party family. Each country is weighted equally, and parties within each country are weighted by their voting share. Details on leader demographics are provided in Table \ref{tab:leader_code_book}}.



\beginQuotable[fig_R1Favorite]
\figureMacro{Decomposition with New Entrants}
{dec_new_entry}
{decomposition/waterfall_entry_46.pdf}
{\textit{Notes:} This figure presents the results of our decomposition exercise, where the residual component is further decomposed into the residuals of new PRRP entrants and the remaining residuals (Equation \ref{eq:residual_decomposition}). We calculate the residual component in Equation \ref{eq:decomposition} separately for PRRP that appeared in our data in the first wave (2005--2009) and those that did not. The black bars present the share of the increase in PRRP support between 2005--2009 and 2017--2020, which we set to 100\% overall. The value for each bar is a weighted average of the decomposition results in each country, with the weights defined as the inverse of the share of PRRP support in the 2017--2020 wave. We use the IVS sampling weights within countries.}
\endQuotable{}



\figureMacro{Average Number of PRRP per Country}
{number_of_prrp}
{manifesto_populist_rise/fr_num_global.pdf}
{\textit{Notes:} This figure shows the average number of PRRP per country that received at least 1\% of the vote. Within each country, the number of PRRP in each year is calculated as the average number of PRRP receiving at least 1\% of the vote in all parliamentary elections in the five-year window centered around that year, based on the CMP data. We then calculate the average number across all \countriesNum{} countries, with all countries weighted equally.}
\clearpage




\clearpage

\beginQuotable[fig_changeLeadership]
\begin{figure}[!htp]
    \centering
    \caption{\normalsize Change in the Residual for Parties With and Without New Leaders}
    \label{fig:newLeaderResid}
    \includegraphics[width=0.8\linewidth]{output/figures/newLeaderResidual/waves46_density.pdf}
    \floatfoot{  \textit{Notes:} This figure shows the distributions of the change in the residual of parties between the first and last waves, as estimated by our model. The figure shows the distribution for parties that kept the same leader and for parties that changed leaders between the two periods.} 
\end{figure}
\endQuotable{}



\clearpage
\beginQuotable[fig_nonPRRPentry]
\figureMacro{Distribution of Counterfactual Vote Shares for New Parties}
{new_party_residual}
{partyResidual/errbar_prrpShrs_vs_all.pdf}
{\textit{Notes:} This figure shows the estimated vote shares that parties first appearing in 2017–2020 would have received if they had existed in 2005–2009. The plot compares new parties classified as PRRP (in red) to all other new parties (Non-PRRP, in black). For each group, the 25th, 50th, and 75th percentiles of simulated vote shares are displayed. These counterfactual shares are estimated by imputing party-specific residuals using a linear model that includes country and party-family fixed effects.}

\endQuotable{}

% Spacing
\setstretch{1.15}


\clearpage

%====================================== APPENDIX TABLES =======================================%



\tableMacro{List of Populist Radical Right Parties}{prrp_parties}
{descriptive_stat/prrp_parties.tex}
{\textit{Notes:} This table provides a list of parties classified as PRRP in our data. Section \ref{subsec:party_class} provides details on how we classify parties.}


\tableMacro{IVS Data}{ivs_data}
{descriptive_stat/dis_stat_table.tex}
{\textit{Notes:} This table provides descriptive statistics on the final dataset analyzed. Each row represents an Integrated Values Survey wave. The observations include only respondents who were successfully matched with the Comparative Manifesto Project data. PRRP support share is the average support for PRRP across the \countriesNum{} countries. \ivsWeights}

 \thispagestyle{empty}
\begin{table}[!htp]
\caption{Decomposition of the Rise in PRRP Support by Country}\label{tab:decomposition_table}
\begin{center}
\centering\begingroup\fontsize{10}{12}\selectfont
\input{output/tables/InferenceTable.tex} 
\endgroup{}
\end{center}
\vspace{-3mm}
\floatfoot{\textit{Notes:} This table presents decomposition results for the 16 countries in our data that had any PRRP support in one of our survey waves (2005--2009, 2011--2013, 2017--2020). Columns (1)-(4) present the counterfactual increase in PRRP support in that country between the 2005--2009 and the 2017--2020 waves, based on Equation \ref{eq:decomposition}. Column (5) sums Columns (1)-(4) to show the overall increase in the share of support for PRRP in that country in our data. The last row presents aggregated results from all countries, weighting each country by the inverse share of support for PRRP in the last wave, and scaling the overall increase to 1, as in Figure \ref{fig:decomposition}. Standard errors are calculated using 200 bootstrap simulations. We use the IVS sampling weights within countries.}
\end{table}
\clearpage


\beginQuotable[Tab_ManifestoRhetoric]
\begin{table}[ht]
    \centering
    \caption{Manifesto Rhetoric Descriptive Statistics}
    \input{output/tables/manifesto_rhetoric_sum_stats.tex}
    \label{tab:manifesto_rhetoric_sum_stats}
    \floatfoot{\textit{Notes:} This table presents descriptive statistics on manifesto style for \finalparty{} election-party pairs. The first two columns present the average of each variable in the 2005--2009 wave and the last two columns present their average in the 2017--2020 wave. The variables used are described in Table \ref{tab:leader_code_book}. More details on data collection appear in Appendix \ref{subsec:style}.}
\end{table}
\endQuotable{}



\beginQuotable[tab_PRRPentry]
%\begin{landscape}
\begin{table}[!htp] 
\centering
\footnotesize
\caption{Populist Radical Right Parties Absent in the 2005--2009 Wave}
\label{tab:party_entry_list}
\input{output/tables/descriptive_stat/party_entry_list.tex}
\vspace{-3mm}
\floatfoot{\textit{Notes:} 
This table provides the full list of PRRP that appear in our data in the 2017--2020 survey wave, but not in the 2005--2009 wave. Column (2) lists the party, column (3) explains why it is missing in the first survey wave, and column (4) provides information on whether alternative radical-right parties existed.}
\end{table}
%\end{landscape}
\endQuotable{}

\clearpage


\beginQuotable[tab_rhetoricCons]
\begin{table}[ht]
    \centering
    \caption{Internal Consistency of Manifesto Rhetoric}
    \input{output/tables/manifesto/manifesto_rhetoric_consistency.tex}
    \label{tab:manifesto_rhetoric_consistency}
    \floatfoot{\textit{Notes:} This table reports split-half correlations for manifesto-level rhetorical measures after demeaning by country-specific averages. For each manifesto, the text is divided into ten equal-sized sections based on textual order. All measures are first computed separately for odd and even sections at the manifesto level, and then demeaned by subtracting the country-specific mean to remove country-level differences. The reported correlation coefficient is the correlation between odd and even sections of the manifestos, calculated across all manifestos. Higher values indicate greater internal consistency. Correlations are computed using complete observations only.}
\end{table}
\endQuotable{}


\tableMacro{IVS Data Matched with CMP}{ivs_data_match}
{descriptive_stat/data_filtering.tex}
{\textit{Notes:} This table describes the Integrated Values Survey data. Row 1 is the total number of respondents in the country-waves we analyze. Row 2 restricts to respondents who supported a specific party. Row 3 restricts to those matched to CMP data from an election within five years of the survey. Row 4 is the final estimation sample, which additionally drops parties supported by fewer than ten respondents.}


\clearpage

\vspace{-2cm}


\beginQuotable[Tab_shy]
\tableMacro{Observed and Predicted Support by Party Family Over Time}{shy_voters}
{shy_voters.tex}
{\textit{Notes:} This table presents the observed party support for respondents who reported it, and the predicted support for respondents who did not answer the relevant questions, for each party family in the 2005--2009 and 2017--2020 waves. The observed support is calculated for people who expressed support in the survey for a specific party, and the predicted support is estimated for people who did not answer either of the questions ``Which political party would you vote for first choice?'' and ``Which political party appeals to you most?''. The prediction model is a LASSO multinomial logit regression estimated separately in each survey wave. IVS sampling weights are used.}
\endQuotable{}


\clearpage


\section{Bliss-Point Model}\label{sec:bliss}
In this section, we discuss the similarities between our model and a simple bliss-point model and explain how we estimate the parameters of the bliss-point model for our robustness analysis in Section \ref{sec:results}.

Assume that voters have a bliss-point vector of party characteristics that is an affine transformation of their observables, $A x_i+b$. Voters support parties that are closest to their bliss point. 
Formally, define the distance between two vectors of party characteristics as
\begin{equation*}
    dist(u,v)^2 = (u-v)'D(u-v),
\end{equation*}
where $D$ is a diagonal matrix with a weakly positive diagonal representing the relative importance of different party characteristics. Using this distance function, we can define the bliss-point utility function as
\begin{equation*}
    U_{ij} = -dist(z_j,Ax_i+b)^2 + \zeta_j + \varepsilon_{ij}.
\end{equation*}

We define $\Phi=2A^T D$, and define $\delta_j = (2b-z_j)'D z_j +\zeta_j$. We can then write the utility function as 
\begin{equation*}
    U_{ij} = x_i \Phi z_j'+\delta_j+ \varepsilon_{ij},
\end{equation*}
which is exactly the utility function we estimate in the first stage (Equation \ref{eq:x*phi*z}). Therefore, our estimation of the first stage will be identical to the case of a bliss-point model. 

This equivalence between the two models implies that any counterfactual change to voter attributes would be invariant to which of the two models we use. Any counterfactual changes in the distribution of the attributes depend only on the first-stage parameters $\Phi, \delta$. Since these parameters would be the same in both models, the counterfactual support for any party would be the same in both models.%\footnote{Our results for voter attributes in Panel C of Appendix Figure \ref{fig:decomposition_robustness} are close but not identical to the results in Figure \ref{fig:decomposition} because of different randomization seeds.} 


By contrast, we need to adjust our estimation of the second stage to accommodate the quadratic terms. Specifically, we rewrite $\delta$ as
\begin{equation*}
    {\delta}_{j,t} = \gamma_{t}z_{j,t}^2 +  \beta_{t} z_{j,t} + \zeta_{j,t},
\end{equation*}
where $z_{j,t}^2=((z_{j,t}^1)^2,...,(z_{j,t}^L)^2)$ is the vector of party characteristics squared. 
Similar to Section \ref{subsec:second-stage}, we take the difference between two consecutive survey waves and get the following equation that extends Equation \ref{eq:oaxaca} to the quadratic case:
\begin{equation*} 
    \Delta^{t+1}_t {\delta}_{j} =  
    \underset{\text{Voter Priorities}}{\underbrace{\Delta^{t+1}_t  \gamma \bar{z_{j}^2}+\Delta^{t+1}_t  \beta \bar{z_{j}}}} +  
    \underset{\text{Party Characteristics}}{\underbrace{\bar{\gamma} \Delta^{t+1}_t z_j^2+\bar{\beta} \Delta^{t+1}_t z_j }}  + 
    \underset{\text{Residual}}{\underbrace{\Delta^{t+1}_t  \zeta_{j}}}.
\end{equation*}

We follow a similar procedure to reduce the dimensionality of the party characteristics vectors, as described in Section \ref{subsec:second-stage} and Appendix \ref{sec:beta_restrcit}. We use the same dimension reduction technique for $z_{j,t}$ and $z_{j,t}^2$: we first reduce the dimension of $z_{j,t}$ to $k$ as before, and then take the squares of the reduced dimension vector. We choose $k=3$, such that, in total, $\Delta^{t+1}_t {\delta}_{j}$ is a linear function of $2k=6$ variables, similar to the five-dimensional $z$ vectors we had in the linear case. 

In Section \ref{sec:results}, we present decomposition results using the bliss-point model described in this appendix. We estimate the parameters $(\Phi_t,\beta_t,\gamma_t)$ as explained here and then decompose the rise of PRRP using Equation \ref{eq:decomposition} as in our main decomposition. The only difference from our original decomposition is that the priority component now includes the parameter vector $\gamma$ as well. 


\section{Data Appendix}\label{sec:data_processing}


\subsection{Supply-side Data}\label{subsec:manifesto_data}
To examine the supply-side characteristics of political parties, we analyze party manifestos, leader demographics, the share of female legislators, the manifesto complexity, and manifesto insult score. Our final dataset comprises \finalparty{} party-election pairs in total. In this section, we discuss how we construct the demographic and rhetoric variables, along with the classification of parties, and party pledges. The full set of variables appears in Table \ref{tab:leader_code_book}.

\subsubsection{Party Classification}\label{subsubsec:classification}
We classify parties into families based on the CMP data. We define parties as center-left if the CMP codes them as social democratic; center-right if they are coded as Christian democratic, or conservative; green if they are coded as ecological; and socialist if they are coded as such. These party families, along with the PRRP, account on average for \haveParfamVoteShare{} of the vote share of all parties coded in the CMP data.



\subsubsection{Demographics of Leaders and Share of Female Legislators}\label{demogrphics}
\beginQuotable[demographicsData]

\paragraph{Party leaders}
We collect demographic information on party leaders  from the leaders' Wikipedia pages. We refer to the Wikipedia page in the original language when available; otherwise, we use the English version. Specifically, we document the leader's gender, years of schooling, age, marital status, number of divorces, and minority background. We define minority as someone who belongs to at least one of the following groups: immigrant, non-Christian, non-white, or non-heterosexual. We also include the log of the average weekly wage in the leader's occupation based on IPUMS 2023 data \citep{flood2024ipums}. 

In cases of party co-chairs, we define the party president as the main leader if both a secretary-general and a president are listed. For parties that enter the election under a collective or dual leadership, we randomly select one individual as the main leader.

\paragraph{Female legislators }
To estimate the share of female legislators for each party in a given year within our dataset, we incorporate data from three external sources: \citet{weeks2023mannerparteien};  \citet{AdamsBrackenGidronHorneOBrienSenk2023}; and the Global Legislator Database \parencite{GlobalLegislatorDatabase2024}. 

The integration process matches each party-year observation in our dataset to corresponding entries in the external sources. We first use \citet{weeks2023mannerparteien}, matching exact party-year observations by CMP ID and assigning unmatched cases to the closest available year for the same party within a five-year window. For unmatched entries, we then apply the same procedure to the dataset used in \citet{AdamsBrackenGidronHorneOBrienSenk2023}, using Parlgov IDs as the merge key. For any remaining unmatched entries, we aggregate legislator-level data from the Global Legislator Database \parencite{GlobalLegislatorDatabase2024} to the party level, calculate the female share as the proportion of female legislators, and again match each party-year observation to the closest available year for the same party within a five-year window. 
%For each party-year entry in our dataset, we incorporated the female share value from the corresponding matched observation in the external sources. 


%which is similar to the data processing approach described in Appendix \ref{subsec:data_processing}.
\endQuotable

\subsubsection{Manifesto Rhetoric: Complexity and Insult Score}\label{subsec:style}


\beginQuotable[styleData]We measure parties' rhetorical complexity and insult score using the text of their manifestos. Whenever available, we rely on the manifesto text provided by the CMP \citep{merz2016manifesto}, which includes manifestos in four formats: text, annotated text, PDF, and annotated PDF. The annotated versions include CMP positions for each clause (quasi-sentence), as segmented by the CMP. Of the \finalparty{} party-election pairs, CMP provides annotated textual data for \cmpparty{} pairs and non-annotated textual data for \cmptext, while \missingparty{} pairs are only available in PDF format or are missing the text. For these \missingparty{} cases, we seek additional sources to supplement our dataset: \capparty{} pairs are retrieved from the Comparative Agendas Project \citep{baumgartner2019cap}, \polidocparty{} from Polidoc \citep{benoit2009polidoc}, and \vaultparty{} from the Manifesto Vault Dataset \citep{dolinsky2024manifestovault}. For the remaining four pairs, we impute the supply-side rhetoric data using random forests.


\paragraph{Complexity}
We employ three measures to capture manifesto complexity. The first measure examines conceptual complexity (i.e., how difficult it is to understand the text). We follow \citet{di2023keep}, who estimate complexity in French and American manifestos. Specifically, we use the Moving Average Type-Token Ratio (MATTR) with a 200-word window \citep{covington2010cutting}. This metric estimates lexical diversity by calculating the ratio of unique words to total words, i.e., the Type-Token Ratio (TTR), and smoothing out text-length effects by taking the moving average. We compute the MATTR after removing stopwords, using standard stopword lists \citep{stopwordsiso, bird2009nltk}. The second measure focuses on syntactic complexity, assessing the average number of words per clause. The CMP segments the manifesto into clauses rather than traditional sentences. Sentences are split into clauses if they contain multiple statements. For the \missingparty{} manifestos not included in the CMP data, we estimate clause counts using country-specific ratios of clauses to sentences derived from CMP data. The third measure is the total word count of the manifesto, which provides a straightforward indicator of document length and so measures complexity at the document level. 
While differences in languages across countries can generate differences in each of these measures, this will not affect our results as long as parties in the same election year write their manifestos in the same language. 

\paragraph{Insult Score}
Our second measure of party rhetoric is an insult score, assessed using the Perspective API \citep{lees2022new}. Specifically, we use the ``insult'' attribute, which measures the degree to which a clause or sentence is perceived as insulting, inflammatory, or negative toward a person or group. The API assigns a score to each clause, where, for instance, a score of 0.9 indicates that 90 out of 100 people consider the clause insulting. Since the API does not support all languages in our dataset, we translate the manifesto content into English for this analysis. The CMP provides DeepL-based translations for all countries except Iceland. For Icelandic manifestos and other manifestos not included in the CMP, we translated them using DeepL and added them to our database. Since the vast majority of manifesto text is not insulting and the differences between very low scores are thus irrelevant for this measure, we focus on the average score of the top 5\% most insulting clauses/sentences in each manifesto.
\endQuotable

\vspace{0.5cm}

\beginQuotable[InternalConsistency]
One concern with rhetoric data is that it can be quite noisy. Therefore, we also test whether these measures are internally consistent. We first split each manifesto into 10 consecutive segments of equal length. We then calculate our rhetoric measures (excluding length of the manifesto) separately for even and odd segments, resulting in two values for each measure-manifesto pair. We demean the results by country-level average (e.g., the insult score for all even segments in a country) to remove cross-country level differences. Finally, we calculate the correlation between even and odd segments in all manifestos. Table \ref{tab:manifesto_rhetoric_consistency} shows that the correlations are relatively high (0.64--0.87) for the three measures we can test, suggesting that our measures capture a meaningful and consistent signal within each manifesto. 
\endQuotable

\subsubsection{Pledges} \label{subsec:pledges}
\beginQuotable[pledgesAppendix] Some voters may ignore general declarations and focus only on concrete pledges that the parties make. To study whether this could affect our results, we repeat our decomposition analysis using only the policy pledges in the manifestos. 

We define pledges according to \citet{thomson2017fulfillment}, where a pledge is characterized as containing language that indicates a commitment to future action or outcomes. This includes both firm commitment language, such as ``we will'', ``we promise to'',  or ``we commit'', and more softly-worded expressions of intent, such as ``we support'' or ``we favor''. To classify all the clauses in the manifestos within our dataset, we use GPT-4o.\footnote{We submitted the following prompt for processing:
For this next assignment, I shall provide you with 50 sentences from a Manifesto, each of them separated with a ``.''.
For each sentence, I want you to analyse if the sentence is a pledge/a commitment. 
For a statement to qualify as a pledge, it must contain language indicating commitment to some future action or outcome. Pledges include both firm commitment language, such as “we will” or “we promise to", ``we commit'', as well as more softly described intention, such as “we support” or “we favor,” as long as parties indicate that they support the action or outcome referred to unequivocally. Please do not indicate that a sentence is a pledge if the platform expresses general statements it wishes would happen instead of concrete policy suggestions. For example, sentences such as ``There should be greater scope for employees, by agreement with the employer, to work more during parts of the year in return for taking longer holidays during other parts of the year.'' Should not be considered a pledge, since it does not specify that it's a party's direct policy. 
On the other hand, sentences like ``We Free Democrats demand a simplification of the documentation requirements for the minimum wage.'' Are clearly a pledge since it indicates an actual policy that the party stands for.   
If the statement is considered a pledge (like I've specified), return 1. If the statement doesn't qualify as a pledge, return 0. Do not provide an explanation, but simply return 0 or 1 according to the instructions. 
Here are the following sentences:}


For this analysis, we rely on the manifesto texts we collected from multiple sources. Since we apply a uniform English prompt, our manifesto content is based on the translated versions, with the translation process consistent with that used in the insult score analysis described in Appendix \ref{subsec:style}. Since this exercise relies on having the \textit{annotated} manifesto text, we could not use the \cmptext{} non-annotated CMP manifestos or the \missingparty{} manifestos that were missing from the CMP. For these manifestos, we impute the pledges and statements using random forests. We predict the share of pledges covering each position for the missing parties, based on this share for parties for which we have data, along with the share of the full manifesto covering each position (which exists for all parties). 

% The agreement rate is here: https://github.com/levyroee/PopulistVoting/issues/282#issuecomment-2450439376
Finally, to test whether our classification procedure is able to distinguish between pledges and declarations, we randomly select 100 clauses and manually classify them as statements or pledges. We reach an agreement rate of 98\%, indicating that the classification is of high quality.
\endQuotable

%I switched the order and started with the supply side to fit the paper. 
\subsection{IVS Data Processing}\label{subsec:data_processing}
This section describes the IVS data processing. We study the following 20 countries:  Austria, Denmark, Estonia, Finland, France, Germany, Greece, Hungary, Iceland, Italy, Lithuania, the Netherlands, Norway, Poland, Portugal, Slovenia, Spain, Sweden, Switzerland, and the UK. All the countries appear in our 2005--2009 and 2017--2020 waves. The following countries also appear in the 2011--2013 intermediate wave: Estonia, Germany, the Netherlands, Poland, Slovenia, Spain, and Sweden.

We analyze over 100 demographic and opinion variables. Since our analysis is data-driven, we do not choose a specific set of variables and instead use virtually all variables that are available for the vast majority of country waves. We process categorical variables in the IVS data to keep the number of potential categories reasonable and merge similar variables when possible. For example, we aggregate the responses to the question about respondents' religion into the following categories: Protestant, Catholic, other Christian, Muslim, Jew, Hindu, Buddhist, and other.

We exclude only five survey questions from our main analysis. Specifically, we exclude a question on general political identity in which respondents are asked to position themselves on a left-right scale. Instead of using it in our main analysis, we use this question as an alternative outcome in Section \ref{subsec:priority_results}.  We also exclude from our data variables that are affected by a party's coalition status: confidence in parliament, confidence in political parties, and confidence in government. While our estimation relies on the assumption that voters choose a party based on their opinions, in some cases the reverse could be true, and the support for a party might shape some of the voters' opinions  \citep{barber2019does}. This is especially common when a new government is elected and, as a result of this change in power, voters immediately change their opinions on the government \citep{hetherington2015washington}. Therefore, we exclude the opinions that tend to change once parties join the governing coalition. To do so, for each opinion in our data, we run a linear regression where the dependent variable is the opinion and the independent variable is whether the party supported by the respondent is part of the governing coalition, as determined by the ParlGov dataset \citep{perlgov}, controlling for party fixed effects and country-wave fixed effects. We exclude from our data the three variables most strongly affected by a party's coalition status. Finally, we exclude one variable that asks directly about priorities. Instead, we estimate priorities separately.\footnote{The survey question we exclude asks about only four specific priorities, which do not seem to be relevant to the priorities we estimated in the data: a stable economy, progress toward a less impersonal and more humane society, progress toward a society in which ideas count more than money, and the fight against crime.} The full list of \demOpNum{} IVS variables, along with the way they were coded, is described in Appendix Table~\ref{tab:ivs_vars}. 

%Appendix Table presents the list of voter attributes included in our dataset and how we process them. 
For all variables, we impute missing values using random forests for each country-wave combination separately. When a value is missing for an entire country-wave combination, we typically exclude the variable from our final dataset. In rare cases where the variable is available for almost all other countries in all survey waves, we impute the values for the specific missing country-wave combination using the nearest survey waves for that country. When the variable is available in both a preceding and a succeeding wave, we impute the variable as a linear interpolation of the mean values in each of these waves, according to the year when each survey was taken. When the variable is available in only a preceding or a succeeding wave, we impute the missing data as the mean value of the available wave. For the imputation process, we also use survey waves conducted before 2005. 

In Appendix Figures \ref{fig:opinion_position_cor} and \ref{fig:heatmap_1_4_subset}, we include earlier survey waves in our analyses. To do so, we removed countries and variables where we did not have sufficient coverage in earlier waves. Specifically, we exclude the following variables from the economic index: \textit{Personal over government responsibility should be increased}; the following variables from the cultural index: \textit{Prostitution justifiable}, \textit{Like idea of army rule}, \textit{Like idea of democracy}, and \textit{Country needs strong leader}; and the following countries: Estonia, Lithuania, Switzerland, Poland, and Norway. 

Our outcome variable is the respondents' preferred party, defined as the party they said they would vote for or the party that appealed to them most. Most surveys asked participants who they would vote for. If participants said they did not know, they were asked which party they supported. In the last EVS wave, participants were only asked which party appealed to them most. We use both questions to define the outcome for as many respondents as possible.

We use the survey weights in all of our analyses with minor modifications. We multiply the weights by a country-specific normalizing constant so that each country has equal aggregate weight. In countries that had two surveys in the same wave (i.e., both EVS and WVS), we multiply the weights for all observations by a survey-specific constant so that both surveys have equal aggregate weight within that country-wave combination. 

%\subsection{Merging Datasets}\label{subsec:data_merging}


\subsection{Linking Party and Voter Data}\label{sec:data-linking}


For our main analysis, we create a novel dataset that links data on voter attributes with data on the characteristics of the parties they support. We first merge the CMP and IVS datasets. These datasets contain data that is key to our model (voters' opinions and demographics, along with parties' positions), and therefore, we only include in our data IVS respondents that we are able to merge with CMP positions. 

We merge the datasets using Party Facts \citep{doring2019party}, a designated platform for linking party datasets. When Party Facts data is unavailable, we manually link parties using their names. By definition, the CMP data is measured around elections, while the IVS surveys are not necessarily conducted close to elections. Therefore, in each survey wave, we assign the party positions from the closest election based on the distance between the mean date when a survey was conducted in a specific country and wave and the date when the election was held. We exclude IVS observations if no CMP data is available five years before or after the survey.\footnote{Although an IVS wave may be composed of both an EVS wave and a WVS wave that were not necessarily conducted in the same year, we assign to each IVS wave and country combination a single date when merging with party positions. We do so in order to assign a single manifesto to each party in each survey wave. However, calculating the mean date separately for the EVS and WVS waves would have only changed the assigned manifesto of a party in a handful of cases and would not have changed the set of observations we are able to match to CMP data within five years.} We also exclude \matchedDropped{} respondents who supported parties supported by fewer than ten respondents in the IVS. We use this threshold to allow for the cross-validation exercise described in Section \ref{subsec:first-stage}. 

Appendix Table \ref{tab:ivs_data_match} shows that we are able to match \matchedPct{} of respondents who expressed support for a party with its CMP data within five years.\footnote{Moreover, \matchedClosestPct{} are matched with a manifesto in the closest election to the survey date.} We do not match all parties for the following reasons: a party may not publish a manifesto, the manifesto of the party may not be coded in the CMP data,\footnote{CMP codes manifestos for parties receiving at least one seat in the elections for the lower house in Western Europe and two seats in elections in Central and Eastern Europe. In some cases, the platforms of parties that met these conditions in the past are also coded.} a party may run in an alliance, and a party may have existed when the survey was conducted but not during the election.

When parties change names or run in various coalitions, it is often not clear if a new party was established or whether the same party ran under a different name or as part of an alliance. We follow the CMP  to deal with this issue and define unique parties according to their CMP ID. 

%The CMP also indicates when one party is a successor of another. However, there are very few pairs of parties where both the predecessor and successor parties appear in our data, and therefore we do not merge predecessor and successor parties. 

After merging the CMP and IVS data, we add the data on leadership demographics. For parties with missing leadership demographics, we use random forests to impute their supply-side characteristics using the supply-side characteristics of all the other parties and the manifesto positions. We impute at least one demographic characteristic for \leaderDemoImputedAnyNum{} parties, impute at least two characteristics for \leaderDemoImputedOverOneNum{} parties (\leaderDemoImputedOverOnePct{}), and impute all demographics for only \leaderDemoImputedAllNum{} parties.

%We impute all demographics for only \leaderDemoImputedAllNum{} parties, impute any demographics for \leaderDemoImputedAnyNum{} parties (\leaderDemoImputedAnyPct{}).


%merge the datasets (CMP and IVS) using Party Facts \citep{doring2019party}, a designated platform for linking party datasets. Where  By definition, the CMP data is measured around elections, while the IVS are not necessarily conducted close to elections.\footnote{The fact that the IVS is not conducted close to elections may imply that IVS respondents are less informed about parties. Moreover, the party they support in the survey may be different from the party they support close to an election. On the other hand, a survey that is not conducted close to an election is less likely to be affected by strategic voting, such as voters taking into account polling results \citep{bursztyn2023identifying}, and thus our sincere voting assumption is more likely to hold in the IVS data.} When merging the datasets, we assign to each party the CMP variables defined for the election closest to when the IVS survey was taken. We include parties in the analysis if at least ten respondents supported them and if we can match them with CMP data five years before or after the year the survey was conducted.\footnote{The CMP typically codes parties that gained at least one or two seats in parliament.} 

%Appendix Table \ref{tab:ivs_data_match} shows that we are able to match \matchedPct{} of respondents who expressed support for a party with its CMP data. This data can be useful for any future research that explores the relations between voter attributes and the characteristics of the parties they support. 


%We merge party data across the various datasets using Party Facts \citep{doring2019party} when possible and manually in other cases. To assign party positions to parties in the IVS data, we first match each party with a party in the CMP data and then, in each survey wave, assign the party positions from the closest election. The closest election is determined based on the distance between the mean date when a survey was conducted in a specific country and wave and the date when the election was held. We exclude IVS observations if no CMP data is available five years before or after the survey.\footnote{Although an IVS wave may be composed of both an EVS wave and a WVS wave that were not necessarily conducted in the same year, we assign to each IVS wave and country combination a single date when merging with party positions. We do so in order to assign a single manifesto to each party in each survey wave. However, calculating the mean date separately for the EVS and WVS waves would have only changed the assigned manifesto of a party in a handful of cases and would not have changed the set of observations we are able to match to CMP data within five years.} We also exclude \matchedDropped{} respondents who support parties supported by fewer than ten respondents in the IVS. We use this threshold to allow for the cross-validation exercise described in Section \ref{subsec:first-stage}. 




\section{Cross-Country Similarities in Voting Decisions} \label{sec:classify_right}

\subsection{Empirical Analysis}

This appendix discusses the assumption that voter priorities are equal across countries. This assumption implies that voters with similar attributes would tend to support parties with similar positions across the continent, since voting weights are a function of attributes and priorities. 
By contrast, if priorities were different across countries, voters with similar opinions could prioritize different topics. For example, in some countries, voters could choose parties mainly based on economic issues, while in other countries, voters with the same demographics and opinions could choose parties mainly based on cultural issues. 

To test this assumption, we regress the support for each party on voter attributes using IVS data. Specifically, for each party $j$, we estimate the linear probability model 
\begin{equation*}
    I_{i,j} = x_{i}'\eta_j+\varepsilon_{ij},
\end{equation*}
where $I_{i,j}$ is a dummy variable for voter $i$ supporting party $j$, using only respondents from the country of that party. 
%We then check how similar are the estimated $\eta_j$ parameters for each pair of parties in two different countries. 
We then test whether parameters that were estimated for a specific party in a specific country can be used to predict voting for a different party in another country. 
We predict the fitted value for each voter and party using the estimated parameters for this party, $\widehat{I_{i,j}}= x_{i}'\widehat{\eta_j}$. 
We then calculate $\rho_{\widehat{I_{i,j_0}},\widehat{I_{i,j_1}}}$, the correlation between the fitted values for each pair of parties with more than fifty supporters, $j_0,j_1$, using all voters in both of the parties' countries. 
%For this exercise, we also include an additional category for other populists, based on the PopuList classification. 

In Section \ref{subsec:link2model}, we show that based on the model in Section \ref{sec:model}, the coefficient $\eta_j$ depends on three things: the voter priorities in the country of party $j$, the party characteristics $z_j$, and the characteristics of the competing parties (to a first-order approximation). Hence, if voter priorities are similar across countries, we expect the correlations within a family to be substantially higher than across families. This is because parties in the same party family tend to hold similar characteristics \citep{de2023party}. However, we note that this exercise is not expected to yield perfect correlations within party families, for at least three reasons. First, there is still heterogeneity in party characteristics even within party families (AfD does not have exactly the same characteristics as UKIP, even though both are PRRP). Second, the probability of voting for a party also depends on the competing parties, which could vary substantially across countries. Third, different countries have different electoral systems, which also affect the party supported.
%\footnote{In addition, in some countries, there are multiple parties in the same category, and the correlation between parties in the same country is smaller, as one cannot support more than one party.}

Appendix Figure \ref{fig:party_correlation} shows the estimated correlation for each party pair.
Red squares indicate similarities between the attributes of the voters of the two parties (positive correlation), while purple indicates dissimilarities. For example, UKIP voters (right column) are particularly similar to AfD voters, such that a model that predicts support for UKIP in the UK can predict support for AfD in Germany.  



We classify parties into families, as explained in Section \ref{subsec:party_class}, and report the average correlations between and within each family. Our classification generates cohesive categories across countries, with positive correlations between parties in the same family. Despite the limitations of using models estimated on completely different countries, we find positive correlations between different PRRP. By contrast, the correlations between PRRP and parties from other families are negative on average. We also find that PRRP are not correlated on average with other (non-right) populist parties, supporting our decision to focus only on right-wing populism  \citep{rooduijn2018unites}. 
These results are fairly consistent with our assumption that voter priorities are equal across countries. 

\subsection{Link to Main Model} \label{subsec:link2model}
To clarify the underlying assumptions of this empirical exercise, we now link it to the model we outline in Section \ref{sec:model}. 
Using Equation \ref{eq:x*phi*z}, the utility of voter $i$ from voting for party $j$ is a linear function of her attributes. 
\begin{equation}\label{eq:single_party}
    U_{i,j}=x_{i}'\Tilde{\eta}_j+\delta_j+\varepsilon_{ij},
\end{equation}
where $\Tilde{\eta_j} = \Phi^{c(j)}_t z_j$, and the matrix $\Phi^{c(j)}_t$ contains the priority parameters of the voters in the country $c(j)$ of party $j$, in period $t$.
Therefore, if two parties hold the same characteristics ($z_{j_0} =z_{j_1}$) and the voters in the countries where these parties are located have the same priorities ($\Phi^{c(j_0)}_t =\Phi^{c(j_1)}_t$),  then the coefficients would be equal $\Tilde{\eta}_{j_0} = \Tilde{\eta}_{j_1}$. 

The probability of supporting a party also depends on the characteristics and residuals of all competing parties in the country ($Z_{t}^{c(j)},\zeta_t^{c(j)}$). Using Equation \ref{eq:P(Z|X)}, we can write the log probability of voter with attributes $x_i$ to support party $j$ as 
\begin{equation*}
    \log P_{i,j} = x_{i}'\Tilde{\eta}_j+\delta_j - C(x_i, (\Phi_{t}^{c(j)},\beta_{t}^{c(j)}),Z_{t}^{c(j)},\zeta_t^{c(j)}).
\end{equation*}
The constant $C$ is the log of the denominator in Equation \ref{eq:P(Z|X)} and guarantees that the probabilities would sum to one. It is common for all parties in a country, and it depends on the voter priorities, as well as the characteristics and residuals of all parties in that country. The coefficient $\eta$ estimated in Equation \ref{eq:single_party} is a linear approximation of this expression. It will be equal for two parties with the same characteristics and same opponents if voter priorities are equal across countries. 

In practice, the party characteristics are not identical across countries, even for parties in the same party family. Moreover, the characteristics of competing parties could be fundamentally different. The electoral systems also differ across countries, which generates differences in the residuals. Therefore, the coefficients in Equation \ref{eq:single_party}  are expected to be more similar within the same family, yet still substantially different. 

\section{Estimation Appendix}\label{sec:beta_restrcit}
In this section, we discuss the dimension reduction in our second-step estimation, which estimates the $\beta$ parameter.
We assume that combinations of party characteristics that generate differences in utility among voters are the same factors that determine the average utility across all voters. 
Formally, the voting weights for every voter are given by Equation \ref{eq:weights}. Using $\Phi=U \Sigma V^T$ to describe the singular value decomposition of matrix $\Phi$, this can be written as $w(x_i) = x_i U \Sigma V^T + \beta$. Defining $\Tilde{\beta} = \beta V$ we can write $w(x) = (x_i U \Sigma +\Tilde{\beta})V^T$. Since we restrict the nuclear norm of $\Phi$ in the first stage, the last components of $x_i U \Sigma$ would be close to zero (assuming the diagonal of $\Sigma$ is ordered). Similarly, we restrict $\Tilde{\beta}$ such that only the first $k$ components are different from zero. Therefore,  $\beta$ has to be a linear combination of the first $k$ components in matrix $V$, such that $\beta \in \text{span}\{[V]_k\}$.\footnote{Using the first $k$ components in an SVD of a matrix yields the best approximation for the matrix for the Frobenius norm based on the Eckart–Young–Mirsky Theorem.} 
%In practice, this is done with linear regression. We calculate $\Tilde{Z} = [V_t]_k]^T Z$ for every IVS wave. We then replace $Z$ with $\Tilde{Z}$ in equation \ref{eq:oaxaca}. 
 
 
 We estimate $\beta$ with a constrained least squares estimator based on Equation \ref{eq:estimating_beta}:
 
 \begin{equation}
 \label{eq:constrained_second_stage}
    \hat{\beta} = \arg \min_{(\beta_t,\beta_{t+2)} \in \text{span}\{[V]_k\}}
    (\Delta^{t+2}_t \hat{\delta}_{j} - 
    \Delta^{t+2}_t  \beta \bar{z_{j}} -  
    \bar{\beta} \Delta^{t+2}_t z_j)^2.
\end{equation}
To solve this equation, we first estimate the following equation using ordinary least squares: 
\begin{equation*}
    \Delta^{t+2}_t {\delta}_{j} =  
     \kappa ([V]_k^T \bar{z_{j}}) +  
    \rho ([V]_k^T \Delta^{t+2}_t z_j)   + 
    \epsilon.
\end{equation*}
The dependent variable is the estimated change in $\delta$ ($\Delta^{t+2}_t \hat{\delta}_{j}$).\footnote{We estimate $\Delta_t^{t+2}\delta_j$ from a pooled regression of $\hat{\delta}{j,t}$ on party-specific linear time trends. We include country-by-wave fixed effects to absorb the arbitrary normalization of $\hat{\delta}{j,t}$ within each option set.} The independent variables are the $k$ linear combinations of the average characteristics ($[V]_k\bar{z_{j}}$) and the change in characteristics ($[V]_k\Delta^{t+2}_t z_j$). We calculate $V$ based on our estimation of $\hat{\Phi}$ in the 2017--2020 wave. We calculate the average party characteristics ($\bar{z_{j}}$) using all survey waves available for this party, and the change in party characteristics ($\Delta^{t+2}_t z_j$) is the difference between the 2017--2020 and 2005--2009 survey waves. We weight each observation by the aggregate vote share of that party in all survey waves. The estimated parameters $\hat{\kappa}, \hat{\rho}$ are two $k$-dimensional vectors. We can then retrieve $\beta$ with a linear transformation.
 \begin{equation*}
 \begin{bmatrix}
 \hat{\beta}_t \\
 \hat{\beta}_{t+2}
 \end{bmatrix} = 
  \begin{bmatrix}
 \hat{\rho}-\frac{\hat{\kappa}}{2}\\
 \hat{\rho}+\frac{\hat{\kappa}}{2}
 \end{bmatrix}[V]_k^T.
\end{equation*}
These values are the solution to Equation \ref{eq:constrained_second_stage}. We choose $k=5$ for our main specification. In Section \ref{sec:results}, we show that other values yield similar results.

\section{Legitimacy}\label{sec:legitimacy}

\beginQuotable[legitimacy]

Particular cultural positions or PRRP may have gained legitimacy in recent decades, and such normalization could result in voters prioritizing these issues. While we cannot test this mechanism directly, we do not find evidence for it based on our survey data. 

First, if support for PRRP was considered illegitimate in the past but gradually became normalized, we might expect to observe underreporting of support for PRRP in our earlier surveys, but not in our later surveys, consistent with the Shy Trump Supporter Hypothesis \citep{coppock2017did}. However, while support in the surveys is slightly lower than the actual vote share of these parties, the difference between survey support and election results does not change between our survey waves. 

Second, some voters may hold opinions that are similar to PRRP positions, but do not consider it legitimate to support these parties, or do not want to express such support, resulting in more voters with opinions similar to PRRP voters among the undecided (voters who do not express support for any party in the survey). In Appendix Table \ref{tab:shy_voters}, we predict which party undecided voters in our survey are expected to support based on their opinions and find that the predicted support for PRRP among undecided voters is similar to its observed support. 

Third, if support for PRRP is becoming more legitimate, one may expect the opinions associated with PRRP to become more socially acceptable as well. However, we show in Section \ref{subsec:attributes} that expressed support for conservative cultural opinions remained stable over time on aggregate.
\endQuotable{}


 \section{Results Appendix: Residual Component}\label{sec:residual_appendix}

In this section, we study the residual component. Subsection \ref{subsec:residual_new} shows that the number of PRRP increased dramatically during the time period studied and that the residual component is driven mostly by these entries. Simple supply and demand explanations for why these parties did not enter earlier are inconsistent with our findings: Subsection \ref{subsec:supply_demand} shows that a simple supply story of PRRP entering the political system cannot explain our results since PRRP existed in the earlier period, even when they did not receive support; Subsection \ref{subsec:res_demand} shows that a simple demand story cannot explain our results because latent demand for these parties existed before they entered. 
Finally, in Subsection  \ref{subsec:residual_alternative}, we propose alternative explanations, including strategic considerations and unobserved support for new PRRP entrants.


\subsection{The Residual Component and New Entries}\label{subsec:residual_new}


 
 \beginQuotable[ResidualEntries]
 Appendix Figure \ref{fig:number_of_prrp} shows that the average number of PRRP increased substantially during the period of our study, suggesting that the entry of new PRRP may drive some of the residual component, as party entry is always attributed to the residual. The figure presents the average number of PRRP that received at least 1\% of the vote across all elections over a five-year window in all the countries in our sample, based on the CMP. 

We next decompose the residual component to measure how much of its rise is explained by new entrants. We calculate the same counterfactual as before, where only the residual component changes over time, and other components are held fixed at their initial levels in 2005--2009, but now separate the overall support for PRRP ($\Pi$) into support for new PRRP entrants ($\Pi_{new}$) and for PRRP that already existed in our data in the 2005--2009 wave ($\Pi_{exist}$). Formally, we decompose the residual component into the following two components: 
\begin{equation}
\label{eq:residual_decomposition}
\begin{array}{ccc}
\underset{\text{New Entrants}}{\underbrace{\int P\left(\Pi_{new}|x_{i};(\Phi_{t},\beta_{t}),Z_{t}^c,\zeta_{t+1}^c\right)f^c_{t}\left(x_{i}\right)\text{d}x_{i}-\int P\left(\Pi_{new}|x_{i};(\Phi_{t},\beta_{t}),Z_{t}^c,\zeta_t^c\right)f^c_{t}\left(x_{i}\right)\text{d}x_{i}}}\\
+ 
\underset{\text{Remaining Residual}}{\underbrace{\int P\left(\Pi_{exist}|x_{i};(\Phi_{t},\beta_{t}),Z_{t}^c,\zeta_{t+1}^c\right)f^c_{t}\left(x_{i}\right)\text{d}x_{i}-\int P\left(\Pi_{exist}|x_{i};(\Phi_{t},\beta_{t}),Z_{t}^c,\zeta_t^c\right)f^c_{t}\left(x_{i}\right)\text{d}x_{i}}}\\
\end{array}
\end{equation} 
New entrants are defined as PRRP that do not appear in our 2005--2009 wave data. Parties typically appear in the CMP data if they enter parliament (or receive at least two seats in  Central and Eastern Europe).

In contrast to the decomposition in Equation \ref{eq:decomposition}, which represents four different counterfactuals, this additional decomposition calculates the impact of the same counterfactual (change in the residuals) on two different sets of parties. 
We calculate the counterfactual rise of new and existing PRRP in each country and aggregate the results as before. The ``New Entrants'' component measures the counterfactual support for new entrants in 2005--2009, if they had the same residual as the later waves, instead of the residual of $-\infty$ that we assign to all parties before they enter.

Appendix Figure \ref{fig:dec_new_entry} shows that the rise of the residual component for new entrants is larger than that of all parties combined. 
%On aggregate, \entryResidualRise{} of the overall rise of PRRP is driven by the residual growth of new parties. 
By contrast, for the existing PRRP, the residual is negative. Hence, based on party characteristics, voter attributes, and voter priorities, we would have expected the growth of these parties to be larger. This could reflect the fact that, due to idiosyncratic reasons, in some countries (e.g., Norway and Switzerland), PRRP gained popularity earlier and therefore might not have had as much potential for additional growth \citep{europe2007populist}. 
\endQuotable{}









\subsection{Supply-Side Explanations}\label{subsec:supply_demand}
\beginQuotable[ResidualSupply]


The rise of the residual due to party entry is unlikely to reflect simple supply factors. 
To study supply-side shocks, we manually examine all cases of PRRP entry in Appendix Table \ref{tab:party_entry_list}.  The table shows that the unexplained rise of new entrants is unlikely to be a simple supply shock in which a new PRRP appeared for the first time, as voters already had the option of supporting a radical right party during the 2005--2009 survey wave. Seven PRRP in seven countries received support in our data in the 2017--2020 wave but not in 2005--2009: Estonian Conservative People's Party (Estonia), Alternative for Germany (Germany), Golden Dawn (Greece), Brothers of Italy (Italy), Forum for Democracy (the Netherlands), Sweden Democrats (Sweden), UK Independence Party (UK). 

In three of these countries (Greece, Sweden, UK), the entering party already existed around 2009, but was not recorded in our data due to insufficient support.\footnote{The Sweden Democrats and UKIP received non-negligible support, as we discuss in the next section.} Hence, these parties ``entered'' the political system in the sense of receiving significantly more votes, but voters could have supported them in earlier elections. 

In two of the remaining countries (Italy, the Netherlands), the new PRRP were established later, but voters had the option to support another PRRP, which had already entered parliament (League in Italy, Party for Freedom in the Netherlands).

Finally, in Germany and Estonia, the current PRRP were formed more recently, but other radical-right parties attempted to enter parliament and failed (National Democratic Party in Germany, Independence Party in Estonia) \citep{europe2007populist}.  

To conclude, in all cases where a PRRP appeared for the first time in our data in 2017--2020, voters had the option to support a PRRP or a similar party in 2005--2009. Therefore, the residual cannot be explained by voters having the option to support a PRRP for the first time. 
\endQuotable{}


\beginQuotable[ChangeLeadership]
Another potential supply-side explanation is that some PRRP entries were driven by a leadership change. Even if PRRP existed in the earlier period, they may have had less capable leaders (e.g., leaders with lower charisma, political acumen, or name recognition). While our data includes information on party leadership, this information focuses on a subset of demographics. 

To test whether we might be missing additional important leader characteristics, we use data on all existing parties where we have estimates for the party's residual and observe the party leader in both periods. We test whether the residual changes more when the party has a new leader. Appendix Figure \ref{fig:newLeaderResid} shows that changes in the residual are similar for parties that kept the same leader between our first and last survey waves and parties whose leadership changed. %Moreover, the residual increases less when leadership changes. 

We note that this analysis suffers from several limitations. First, the timing of the leader replacement is not random, and we cannot interpret these estimates as the causal effect of leadership. Second, the analysis is based on all existing parties, while new leadership may matter more for new parties or for PRRP. Still, based on this analysis, we do not find evidence that new party leadership tends to change voters' support. 
\endQuotable{}


\beginQuotable[ResidualSupply2]
We cannot conclusively rule out alternative supply-side developments. First, the parties that entered may have been substantially different from the parties previously available to voters (or from the previous version of the same party before it received support). 
%For example, the Estonian Conservative People’s Party may have had more organizational capacity than the Estonian Independence Party, which could have allowed it to increase support. Similarly, 
For example, UKIP may have professionalized over time, allowing it to attract more voters \citep{goodwin2015ukip}. 
%Finally, Brothers of Italy has more pro-welfare positions than the League, and these positions could explain some of the support it received. 


Second, PRRP may have benefited from changes in the party characteristics that are not captured in our data. While we attempted to capture party positions, rhetoric, and leadership demographics, other changes could have occurred. For example, several PRRP have made concerted efforts to detoxify their electoral brand. The French National Front was renamed Rassemblement National in 2018 \emph{``in an attempt to rebrand it as a more acceptable mainstream party with a wider social agenda''} \citep[393]{moffitt2022mainstream}. If this change was not captured by changes in the insult score of the party platform, we may not capture it in our data. Another supply-side strategy involves recruiting better candidates for office or providing candidates with stronger political training \citep{de2011fortuyn}. This strategy does not operate via party positions but rather via the ability to convey these positions more persuasively to the public. We would not attribute the better candidates to a change in party characteristics if this change is not captured by, or correlated with, the demographics of the party leadership, the manifesto rhetoric, or the party positions. 

Furthermore, we cannot entirely dismiss supply-side theories that are outside the scope of our model. One such hypothesis posits a significant delay between changes in party positions and voter responses. Scholars of electoral politics have argued that the ideological convergence of mainstream parties during the 1990s opened a space for PRRP to emerge \citep{berman2019populism}. Since this shift in positions occurred long before the rise of PRRP support, it is not captured in the party characteristics component of the decomposition analysis. If ideological convergence in the 1990s affected voting decisions in the 2010s, the effect would be attributed to the residual component. 
\endQuotable{}

\subsection{Demand-Side Explanations}\label{subsec:res_demand}

\beginQuotable[ResidualDemand]

The counterfactual analysis in  Section \ref{subsec:residual_new} suggests that demand for PRRP that entered in the later survey wave already existed in 2005--2009. 
Specifically, the party entry residual component measures the latent demand for new PRRP entrants in the 2005--2009 wave. This component calculates the counterfactual support for new PRRP entrants if they had entered before, holding voter attributes, voter priorities, and party characteristics fixed at their values in the 2005--2009 wave, and using the residuals of the 2017--2020 wave (and the 2017--2020 party characteristics for new entrants). We find that if these parties had entered earlier, their support would have reached more than 70\% of the aggregate weighted PRRP rise throughout the entire period. Still, these parties failed to enter, implying that an increase in voter demand is unlikely to be the primary driver behind PRRP entry. 

A possible concern with using the 2017--2020 residuals is that they may reflect a large demand for PRRP that only occurred in later waves. To address this concern, we apply an alternative method to impute the residuals based on party family, and use them to measure latent demand (the counterfactual support for all parties in the 2017--2020 wave, if they had contested in 2005--2009). We first regress the estimated residuals in the 2005--2009 wave for parties that already existed on indicators for country and party family. We then use this model to predict residuals for new entrants. The imputed residuals allow us to calculate the latent demand for the new entrants: we estimate the counterfactual support by fixing all four components at their 2005--2009 values,  while using the imputed residuals based on party family for new entrants (along with the party characteristics from the 2017--2020 wave, as before). 
 
 The upper boxplot in Appendix Figure \ref{fig:new_party_residual} shows the 25th, 50th, and 75th percentiles of counterfactual support for new PRRP entrants. We find large latent demand for new PRRP, similar to results using the residuals from the 2017--2020 wave. This implies that PRRP entrants already had latent demand in 2005--2009.

  \endQuotable{}

\subsection{Alternative Explanations} \label{subsec:residual_alternative}

\beginQuotable[ResidualAlternative]
In this subsection, we first discuss two alternative explanations for the residual component: strategic considerations and unobserved prior PRRP support. We then present suggestive evidence supporting these mechanisms. 

\paragraph{Strategic considerations}
One potential explanation for the residual increase due to new entrants is strategic considerations and coordination failures. Many voters coalesce around a party only when they anticipate that the party will have substantial support. This is because they think larger parties have more influence, because of a bandwagon effect, or because they are afraid that smaller parties would not pass the election threshold \citep{freden2014threshold}. Indeed, PRRP have found it easier to enter parliament in countries with a lower electoral threshold \citep{norris2005radical}.


Therefore, one plausible hypothesis is that some voters did not support PRRP in the earlier survey wave since these parties lacked sufficient support. The prioritization of cultural issues may have generated a critical mass of support for PRRP. This critical mass solved the coordination problem and therefore increased the support for new PRRP entrants, even beyond the direct expected increase from the change in priorities. Such excessive support would be attributed to the residual component, even if it was triggered by a change in priorities. 

\paragraph{Unobserved prior PRRP support}

Another reason for the rise in the residual component, driven by new entries, is unobserved support for PRRP in the 2005--2009 wave. This occurs due to two types of data constraints. First, some of the new entrants already existed in the 2005--2009 wave but were not included as an option in the IVS questionnaire. Therefore, we underestimate their true level of support at that time. The most extreme example is the Sweden Democrats. The party received close to 3\% of the vote in the 2006 elections, and its support increased to around 6\% by 2010. However, it was not asked about explicitly in the IVS, and therefore its measured support is mechanically set to 0\% in our data.
In the UK, UKIP and BNP received combined support of approximately 5\% in surveys conducted in 2009 \citep{ipsos2009aggregate} and more than 5\% of the vote in 2010, but neither party was included in the 2009 IVS questionnaire. Besides these two examples, we have not found other cases of prominent PRRP that were not asked about in the surveys.

Second, splits and mergers could also generate a rise in the residual component through new entrants due to changes in party classification. 
For example, in Italy, the Brothers of Italy split from Berlusconi's Forza Italia. Since Brothers of Italy is classified as a PRRP and Forza Italia is not, this generates an increase in support for a new PRRP entrant, which is attributed to the residual.
In Estonia, the Conservative People's Party was formed when the People's Union merged with a nationalist movement. While the new party is classified as PRRP, the People's Union is not. We classify this as a party ``entry'' in our data and attribute it to the residual component. However, part of the support for this party reflects the preexisting support for the previous party, and not a rise in support for PRRP.  In both of these cases, the unobserved support for PRRP before they appeared in our data would result in a mechanical increase of the residual component.

While unobserved prior PRRP support artificially inflates the measured growth in PRRP support and, in particular, the residual component, this does not affect our main conclusion. Since all the additional growth is attributed to the residual, unobserved prior support does not change the \emph{relative} importance of our three main components (voter priorities, voter attributes, and party characteristics). However, it still makes the residual larger at the expense of these components.

\subsubsection{Comparing Latent Demand for PRRP and Non-PRRP}\label{subsubsec:nonPRRP}
%Shin, make sure this is clear to you. 

Strategic considerations and unobserved prior support both imply that demand for PRRP existed before their first appearance in the data, consistent with the results in Section \ref{subsec:res_demand}.
Moreover, these two factors are not unique to PRRP and hence should predict similar latent demand for PRRP and non-PRRP entrants. By contrast, supply and demand explanations would not predict latent support for all parties, regardless of whether they are PRRP. For example, if voters wanted to support PRRP but these parties did not exist, we would expect to observe substantially greater latent demand for PRRP entrants compared to non-PRRP entrants.

To test whether latent demand was unique to PRRP entrants, we measure the latent demand for non-PRRP as well. We use the same counterfactual described in Section \ref{subsec:res_demand}, where we impute the residual for new entrants based on party family. 
The lower boxplot in Appendix Figure \ref{fig:new_party_residual} shows the 25th, 50th, and 75th percentiles of the counterfactual support for new entrants and indicates similar latent demand for non-PRRP entrants. This similarity between PRRP and non-PRRP is consistent with strategic considerations and unobserved prior support, and less consistent with supply and demand shocks that are unique to PRRP.  

\endQuotable{}



\clearpage


%also exists for other countries. unobserved factors, other things we do not model. 
%The decline in residual for existing parties captures a decrease in support for parties that is not driven by voter attributes, voter priorities, or party positions. This could be because in some countries such as Norway or Switzerland, PRRP already had large support in the 2005--2009 survey wave that was driven by idiosyncratic shocks.  

%We further find that the average share of countries with at least one populist party increased from \populistSharePre{} to \populistSharePost{}. 
%The model confirms that the entry of populist parties contributes to their increased support. 




%\subsection{Additional results}
%[can have a section with other counterfactual, can characterize marginal voters and show where they came from]

%%%% Appendix B Tables %%%%%%


\renewcommand{\thefigure}{B.\arabic{figure}}
\setcounter{figure}{0}     
\renewcommand{\thetable}{B.\arabic{table}}
\setcounter{table}{0}  

%\begin{landscape}

\section{Supplementary Data Tables}


% Manifesto Table
\input{output/tables/manifesto/manifestoVariablesExtended.tex}
\vspace{-4mm}
{\noindent \footnotesize{\textit{Notes:} This table presents the descriptions of all major (three-digit) CMP positions. All of these positions are included in the decomposition exercise. The third and fourth columns show which positions are part of the CMP economic and cultural indices (the culture index is termed ``Society'' in the CMP). These are the primary party position indices used in the paper. The fifth column shows which positions are part of the European integration index (we multiply this index by $-1$ so larger values are more conservative). The next four columns refer to alternative constructions of economic and social indices using CMP positions. These indices are only used in Figure \ref{fig:alternative_cmp}. The next column shows the components of the expanded Hawk-Dove index. The original Hawk-Dove index is proposed by \citet{kiratli2021hawkdove}. We expand it by subtracting the dovish positions from the hawkish ones. For example, we take into account `Military: Positive' in addition to (the negative value of) `Military: Negative'. The last column shows the positions included in the government efficiency index. The expanded Hawk-Dove, European integration, and government efficiency indices are only used in Figure \ref{fig:alternative_indices}.
%Prosser refers to the economic and social indices proposed in \cite{prosser2014building}, while Bakker and Hobolt refers to the economic and libertarian-authoritarian indices proposed in \cite{bakker2013measuring}.
}}

%\end{landscape}

\clearpage

\input{output/tables/party_leaders_code.tex}
\vspace{-4mm}
{\noindent \footnotesize{\textit{Notes:} This table presents descriptions of the leadership demographics and rhetoric variables for \finalparty{} party-election pairs in our final analysis data.}}

\clearpage

% IVS Table
\input{output/tables/ivsVariables.tex}
\vspace{-4mm}
{\noindent \footnotesize{\textit{Notes:} This table presents all IVS variables included in our data. The third column provides information on how variables were processed. The fourth column shows which variables are part of the economic and culture opinion indices used in Figure \ref{fig:heat_map} and Appendix Figure \ref{fig:opinion_position_cor}.}}

\clearpage

\printbibliography

\end{refsection}


% For Grammarly
\if\grammarlyVersion1
\theendnotes
\fi

% Try adding subfigure to A.6:

\end{document}